\documentclass[11pt]{article}
\usepackage{amsmath,amssymb,amsthm,mathtools,mathrsfs,hyperref}
\newtheorem{theorem}{Theorem}[section]
\newtheorem{proposition}[theorem]{Proposition}
\newtheorem{lemma}[theorem]{Lemma}
\newcommand{\C}{\mathbb C}\newcommand{\R}{\mathbb R}
\newcommand{\B}{\mathscr B}\newcommand{\M}{\mathcal M}
\title{Exact receivers from earlier Split Zero mathematics}
\date{19 September 2026}
\begin{document}\maketitle
This mathematical intake recovers several distinct constructions from
earlier manuscripts and proves their maps into the theta-source and
full-multiplicity zero objects used in the RH programme. It does not evaluate
the currently unresolved seven absolute allocations or projected signs.
All source versions and precise proof locators are recorded in the accompanying
\texttt{SOURCE\_RECEIVERS.md}. The statements below are proved here rather than
inferred from names or similarities of constructions.

\section{The original theta source from the earlier circle heat construction}
Retain the physical variable $x>0$, $D=-x\partial_x$, and
\[
 \phi_*(x)=(4\pi^2x^4-6\pi x^2)e^{-\pi x^2},\quad
 \Theta\phi(x)=2\sum_{n\geq1}\phi(nx),\quad f_r=D^r\Theta\phi_*.
\]
The source space is
\[
 V=\{\phi\in\mathcal S(\R):\phi(-x)=\phi(x),\ \phi(0)=0,
                  \ \int_\R\phi=0\},\qquad \mathcal H=L^2(\R_{>0},dx).
\]
These are the original conventions in (TC1)--(TC6) of the theta-chain text.
The older corrected spline manuscript uses the circle coordinate $z\in\R/\mathbb Z$,
\[
 h_0(z)=\mathbf1_{[0,1/2)}(z)-\mathbf1_{[1/2,1)}(z),\qquad
 N=-\frac1{4\pi^2}\partial_z^2,\qquad u_0(z,t)=e^{-tN}h_0(z).
\]
Set $C(t)=\partial_z u_0(0,t)$ for $t>0$.

\begin{theorem}[Exact dyadic source and its full tail]\label{thm:dyadic}
For each integer $J\geq0$ and $r\geq0$, put $s_J=2^{J+1}$ and
\[
 f_r^{[J]}(x)=\frac14D^{r+1}(D-1)
                 \sum_{j=0}^{J}C(\pi4^j x^2).
\]
Then, with no rescaling of $x$ in the receiving theta complex,
\begin{align}
 C(t)&=4\sum_{\substack{k\in\mathbb Z\\k\text{ odd}}}e^{-tk^2},\label{eq:C}\\
 f_r(x)-f_r^{[J]}(x)&=f_r(s_Jx),\label{eq:exact-tail}\\
 \|f_r-f_r^{[J]}\|_{L^2(dx)}&=s_J^{-1/2}\|f_r\|_{L^2(dx)}.\label{eq:tail-norm}
\end{align}
Moreover $f_r^{[J]}=\Theta(\phi_r-\phi_r(s_J\,\cdot))$ with
$\phi_r=D^r\phi_*\in V$. Thus every truncation lies in the original
theta image, and its complete Mellin transform is
\[
 \M f_r^{[J]}(w)=(1-s_J^{-w})w^r\,2\xi(w).
 \tag{D1}
\]
\end{theorem}
\begin{proof}
For $k\ne0$, direct integration gives
$\widehat h_0(k)=(1-(-1)^k)/(\pi i k)$ and $\widehat h_0(0)=0$.
The heat Fourier series and all its spatial derivatives converge absolutely
for $t>0$, giving \eqref{eq:C}. With
$\vartheta(v)=\sum_{n\in\mathbb Z}e^{-\pi n^2v}$, partitioning the
integers into odd and even ones gives
\[
 \tfrac14 C(\pi4^j x^2)=\vartheta(4^jx^2)-\vartheta(4^{j+1}x^2).
\]
The finite sum therefore equals
$\vartheta(x^2)-\vartheta(s_J^2x^2)$ exactly. Also
$D(D-1)e^{-\pi x^2}=\phi_*(x)$; termwise differentiation of the
Gaussian series on every compact subset of $(0,\infty)$ gives
$D^{r+1}(D-1)(\vartheta(x^2)-1)=f_r(x)$.
Since $D$ commutes with $T_sF(x)=F(sx)$, this proves
\eqref{eq:exact-tail}. The Gaussian source and its $D$-iterates have the
two required zero moments. Poisson summation gives rapid decay of their
theta sums at zero as well as infinity, so they belong to $\mathcal H$.
Changing variables $y=s_Jx$ proves \eqref{eq:tail-norm}.

Dilation preserves evenness, the value zero at zero and the zero integral
of $V$. The sum defining $\Theta$ commutes with dilation, proving the
displayed source identity. For $\Re w>1$, termwise Mellin integration and
two integrations by parts give
\[
 \M f_0(w)=w(w-1)\pi^{-w/2}\Gamma(w/2)\zeta(w)=2\xi(w).
\]
Both sides extend to entire functions by the rapid decay just proved;
the convention for $\xi$ agrees with Apostol's formula recorded in
\cite{DLMF25}. Mellin transformation of $D$ and $T_s$ multiplies by $w$
and $s^{-w}$ respectively. This proves (D1).
\end{proof}

\begin{proposition}[Complete Gram receiver]\label{prop:dyadic-gram}
Let $E$ be any one of the original finite packet coordinate spaces,
$F_L:\C^{L+1}\to\mathcal H$, $F_Le_r=f_r$, and
$C_E:E\to\C^{L+1}$ any fixed original coefficient map. Set
$P=F_LC_E$, $P_J=(1-T_{s_J})P$, and retain any original
$R_{\rm ref}:E\to\mathcal H$. For $R=R_{\rm ref}+P$ and
$R_J=R_{\rm ref}+P_J$ one has
\begin{align}
 \Delta_J&=T_{s_J}P, & \|\Delta_J\|&=s_J^{-1/2}\|P\|,\nonumber\\
 R_J^*R_J-R^*R
   &=-R^*\Delta_J-\Delta_J^*R+\Delta_J^*\Delta_J,\tag{D2}\\
 \|R_J^*R_J-R^*R\|
   &\leq2s_J^{-1/2}\|R\|\|P\|+s_J^{-1}\|P\|^2.\tag{D3}
\end{align}
The same exact dilation norm holds for the Hilbert--Schmidt norm.
If $q:\B\to\B/\Theta V$ is the original quotient, then
$qR_J=qR$. All their full Mellin jets at actual zeros of $\xi$ agree.
\end{proposition}
\begin{proof}
$T_s^*T_s=s^{-1}I$ on $L^2(dx)$, which proves both norm equalities.
Expand $(R-\Delta_J)^*(R-\Delta_J)$ without discarding either cross term
to obtain (D2); the triangle and submultiplicative norm inequalities
give (D3). Every column of $\Delta_J$ belongs to $\Theta V$ by the
preceding theorem. Its Mellin transform is a polynomial times
$s_J^{-w}2\xi(w)$, so it vanishes to every original zero multiplicity.
\end{proof}

The exact receiving map is
$\phi_r\mapsto\phi_r-\phi_r(s_J\,\cdot)$ in $V$.
It does not preserve the finite polynomial--Gaussian span
$\operatorname{span}(\phi_0,\ldots,\phi_L)$: the second term has a
different Gaussian scale. Consequently (D2) compares explicitly specified
representatives of the same original theta class; it does not identify two
finite-cutoff attained minima. The useful gain for the programme is a
geometric heat realization of the full original source, with an exact error
in the physical Hilbert norm and with all Mellin multiplicities preserved.

\section{Full-multiplicity observations on the arithmetic-site approximation spaces}
The earlier arithmetic-curve manuscript has the spaces
\[
 H_{\rm B}=L^2((1,\infty),y^{-2}dy),\quad
 F_\alpha(y)=\alpha\lfloor y\rfloor-\lfloor\alpha y\rfloor
            =\{\alpha y\}-\alpha\{y\},\quad 0<\alpha\leq1,
\]
\[
 V_Q=\operatorname{span}\{F_{a/Q}:1\leq a<Q\},\qquad
 x_Q=\operatorname{dist}_{H_{\rm B}}(1,V_Q)^2.
\]
We retain this coordinate and measure; no identification with the theta
Hilbert metric is implicit. This construction imports the human
Nyman--Beurling and B\'aez-Duarte criterion, cited explicitly in
\cite{Beurling,BD}. The new finite receiver below is proved independently
of that equivalence.

\begin{lemma}[Exact Mellin observation]\label{lem:NB-mellin}
For $\Re w>0$,
\[
 \int_1^\infty F_\alpha(y)y^{-w-1}\,dy
       =\zeta(w)\frac{\alpha-\alpha^w}{w}.\tag{B1}
\]
The apparent singularity at $w=1$ on the right is removable.
\end{lemma}
\begin{proof}
$F_\alpha$ is bounded, so the left side is holomorphic on $\Re w>0$.
For $\Re w>1$, write
$\lfloor y\rfloor=\sum_{n\geq1}\mathbf1_{[n,\infty)}(y)$ and
$\lfloor\alpha y\rfloor=\sum_{n\geq1}\mathbf1_{[n/\alpha,\infty)}(y)$.
Absolute convergence justifies integration, yielding
$\alpha\zeta(w)/w-\alpha^w\zeta(w)/w$. The numerator vanishes at $1$,
so analytic continuation proves (B1) on the stated half-plane.
\end{proof}

\begin{theorem}[Confluent Gram lower bound with every multiplicity]\label{thm:NB-jets}
Let $Z_+$ be a nonempty finite set of actual zeta zeros
$\rho$ with $\Re\rho>1/2$, retaining their full multiplicities $m_\rho$.
Use the original local coordinates $(w-\rho)^j$, $0\leq j<m_\rho$,
of $E_{Z_+}=\prod_{\rho\in Z_+}\C[w]/(w-\rho)^{m_\rho}$.
Define
\[
 (J_+f)_{\rho,j}=\frac1{j!}\int_1^\infty
            f(y)(-\log y)^j y^{-\rho-1}\,dy,
 \quad b_{\rho,j}=\frac{(-1)^j}{\rho^{j+1}}.
\]
Then $J_+:H_{\rm B}\to E_{Z_+}$ is a bounded surjection,
$J_+F_\alpha=0$, $J_+1=b$, and its positive definite Gram matrix is
\[
 \mathsf G_{(\rho,j),(\sigma,k)}
  =\frac{(-1)^{j+k}(j+k)!}
     {j!k!(\rho+\bar\sigma-1)^{j+k+1}}.\tag{B2}
\]
For every $Q$,
\[
 x_Q\ \geq\ b^*\mathsf G^{-1}b\ >0.\tag{B3}
\]
In particular one simple observation at $\rho=\beta+i\gamma$ gives
$x_Q\geq(2\beta-1)/|\rho|^2$.
\end{theorem}
\begin{proof}
With inner products conjugate-linear in their first variable, the
representer of the $(\rho,j)$ observation is
$k_{\rho,j}(y)=(-\log y)^j y^{1-\bar\rho}/j!$.
Its squared norm is finite because $2\Re\rho-1>0$.
Thus $J_+^*v=\sum k_{\rho,j}v_{\rho,j}$.
Setting $t=\log y$ and integrating $t^{j+k}e^{-at}$ gives (B2), with
$a=\rho+\bar\sigma-1$ of positive real part.
The representers are independent: a linear relation is an
exponential polynomial $\sum_\rho p_\rho(t)e^{(1-\bar\rho)t}=0$;
successively applying the differential operators
$(\partial_t-(1-\bar\sigma))^{m_\sigma}$ for all but one exponent
isolates a nonzero invertible operator on its polynomial coefficient.
Every coefficient therefore vanishes. This proves $\mathsf G>0$ and
$J_+J_+^*=\mathsf G$, hence surjectivity.
Differentiating (B1) through order $m_\rho-1$ proves $J_+F_\alpha=0$;
all differentiation is justified by the logarithmic integrable bounds.
Direct integration of $1$ gives $b$.

The exact minimum-norm lift of $b$ is
$S_+b=J_+^*\mathsf G^{-1}b$.
Every $e$ with $J_+e=b$ decomposes orthogonally as
$S_+b+(e-S_+b)$, with the second summand in $\ker J_+$.
Consequently $\|e\|^2\geq b^*\mathsf G^{-1}b$.
Apply this to $e=1-v$, $v\in V_Q$, and take the minimum. Since the
coordinate $b_{\rho,0}=1/\rho$ is nonzero, the lower bound is positive.
For one value observation $\mathsf G=(2\beta-1)^{-1}$, giving the last
formula.
\end{proof}

\begin{proposition}[The Beurling observation enters the original theta class]
For the same actual finite divisor $Z_+$, use its original reference
section $R_{Z_+}$ of (T2), with the unit $j_{h_{Z_+}}(h_{Z_+}/(2\xi))$.
Then
\[
 \mathcal T_+=R_{Z_+}J_+:H_{\rm B}\longrightarrow\B
 \tag{B4}
\]
is continuous into the original test-function topology and satisfies
\[
 J_{\rm all}\mathcal T_+=\iota_+J_+,\qquad
 q\mathcal T_+=\sigma_{Z_+}J_+,\qquad
 \mathcal T_+F_\alpha=0,\qquad
 q\mathcal T_+1=\sigma_{Z_+}b\ne0.
 \tag{B5}
\]
Here $\iota_+$ is full-jet extension by zero, and all other maps are
the original maps of the theta complex. Thus (B3) is the exact attained
quotient metric of a second analytic presentation of these same finite
theta classes.
\end{proposition}
\begin{proof}
$J_+$ is bounded and has finite-dimensional codomain. A linear map
from that finite-dimensional space into $\B$ is continuous; concretely
every defining seminorm of $R_{Z_+}v$ is bounded by a finite sum of the
seminorms of its fixed columns times the coordinates of $v$.
The identities follow from Theorem~\ref{thm:NB-jets} and (T2), including
the original complete local unit. Injectivity of $\sigma_{Z_+}$ and
$b\ne0$ prove the final nonvanishing. The quotient norm on
$H_{\rm B}/\ker J_+$ is exactly $v^*\mathsf G^{-1}v$ by the
minimum-norm lift proved above. This assertion specifies that quotient
metric; it does not equate it with any theta $L^2$ metric.
\end{proof}

\begin{proposition}[Exact arithmetic-site tower receiver]\label{prop:tower}
Put $\Lambda_Q=\operatorname{lcm}(1,\ldots,Q)$ and
$R_Q=\operatorname{span}\{F_{1/n}:2\leq n\leq Q\}$.
Then $R_Q\subset V_{\Lambda_Q}$ and, for
$d_Q=\operatorname{dist}(1,R_Q)$, $x_{\Lambda_Q}\leq d_Q^2$.
For $Q_0\geq3$, $Q_{j+1}=\Lambda_{Q_j}$ is strictly increasing and
every positive integer divides some $Q_j$.
Every off-line actual zero supplies the fixed positive lower bound (B3)
at all these levels, using its right-half-plane functional-equation partner
when necessary.
\end{proposition}
\begin{proof}
For $n\leq Q$, $1/n=(\Lambda_Q/n)/\Lambda_Q$, which gives the
inclusion and distance inequality. Because $Q$ and $Q-1$ are coprime,
$\Lambda_Q\geq Q(Q-1)>Q$ for $Q\geq3$. Hence the tower grows without
bound and once $Q_j\geq n$, $n\mid Q_{j+1}$. If $\Re\rho<1/2$,
$1-\rho$ has real part greater than $1/2$ and the same full
multiplicity, by the functional equation \cite{DLMF25}. Apply the preceding
theorem at that actual zero; no hypothetical replacement spectrum is used.
\end{proof}

For clarity, the relation to the usual human criterion is exact. Under the
isometry $Uf(x)=f(1/x)$ from $H_{\rm B}$ to $L^2(0,1)$,
$UF_{1/n}=\rho_n-n^{-1}\rho_1$, where
$\rho_n(x)=\{1/(nx)\}$. Extend this difference by zero to $(1,\infty)$.
If $g=\sum c_n\rho_n$ on $(0,\infty)$ and $A=\sum c_n/n$, then
$g(x)=A/x$ for $x>1$, and
$g-A\rho_1=\sum c_n(\rho_n-n^{-1}\rho_1)$ is supported in $(0,1)$.
For $\chi=\mathbf1_{(0,1)}$,
$|A|=\|g-\chi\|_{L^2(1,\infty)}$, so
\[
 \|g-A\rho_1-\chi\|_{L^2(0,\infty)}
 \leq (1+\|\rho_1\|_2)\|g-\chi\|_2.
\]
Thus the compensated spaces have exactly the same closure criterion as
the integer spaces in B\'aez-Duarte's theorem. Together with the finite
divisibility inclusion above this proves that decay $x_{Q_j}\to0$ on the
cofinal tower is precisely this established RH criterion. No decay estimate
has been assumed or proved here. What is gained is an exact finite,
full-multiplicity observation and its computable positive obstruction at
every arithmetic-site stage. In the arithmetic-curve construction the
framed groups $(1/Q)\mathbb Z$ and monoids $(1/Q)\mathbb N$ carry these
same stages; the transition $Q\mid M$ becomes the monoid Frobenius
$n\mapsto (M/Q)n$ in their framed integral coordinates.

\section{The old compactified jet tail receives the actual theta quotient}
Retain the original spaces
\[
 \B=\{F\in C^\infty(\R_{>0}):\sup_{x>0}x^b|D^rF(x)|<\infty
                      \text{ for every }b\in\mathbb Z,r\geq0\},
 \qquad Q=\B/\Theta V.
\]
Let $Z_\infty$ be the discrete set of all actual nontrivial zeros,
$E_\rho=\C[w]/(w-\rho)^{m_\rho}$, and
\[
 E_{\rm fin}=\bigoplus_{\rho\in Z_\infty}E_\rho,\quad
 E_{\rm all}=\prod_{\rho\in Z_\infty}E_\rho,\quad
 T=E_{\rm all}/E_{\rm fin}.
\]
The earlier packet-geometry manuscript constructs exactly this last
tail as the stalk at infinity of its compactified product-of-jets sheaf.
The following receiver retains both the global tail and the analytic
kernel invisible to every zero jet.

\begin{theorem}[Global full-jet and tail receiver]\label{thm:global-tail}
The map
\[
 J_{\rm all}:\B\longrightarrow E_{\rm all},\qquad
 F\longmapsto\bigl(j_\rho\M F\bigr)_\rho
\]
annihilates $\Theta V$ and descends to a $\C[w]$-module map
$\overline J:Q\to E_{\rm all}$, with $w$ acting by $D$ on $Q$.
The original finite reference sections induce a canonical injective
$\C[w]$-module map $\sigma:E_{\rm fin}\to Q$ satisfying
$\overline J\sigma=\iota$, the finite-support inclusion.
Writing $K=\ker\overline J$ and
$T_{\rm arith}=\operatorname{im}(Q\xrightarrow{\overline J}
E_{\rm all}\to T)$, there is an exact sequence
\[
 0\longrightarrow K\longrightarrow Q/\sigma(E_{\rm fin})
       \longrightarrow T_{\rm arith}\longrightarrow0,
 \qquad T_{\rm arith}\hookrightarrow T.\tag{T1}
\]
The space $T_{\rm arith}$ is nonzero.
\end{theorem}
\begin{proof}
Every Mellin transform of $F\in\B$ is entire and can be differentiated
under its integral, by the bounds defining $\B$ at both endpoints.
For $\phi\in V$, the original Mellin identity is
$\M\Theta\phi(w)=2\zeta(w)\M\phi(w)$ for $\Re w>1$.
The function $\M\phi$ is holomorphic for $\Re w>-2$, since $\phi$
is even and vanishes to order at least two at zero. It follows by
continuation that the displayed product vanishes to every full
nontrivial-zero multiplicity. Thus $J_{\rm all}\Theta=0$.
Integration by parts gives $\M DF(w)=w\M F(w)$ with vanishing endpoint
terms. This proves the module assertion.

For a finite subset $Z$, let
$h_Z(w)=\prod_{\rho\in Z}(w-\rho)^{m_\rho}$, $d_Z=\deg h_Z$, and
$g(w)=2\xi(w)$. If $u\in E_Z=\C[w]/h_Z$, let $r_Z(u)$ be the
unique polynomial of degree $<d_Z$ representing $j_{h_Z}(h_Z/g)u$.
The original reference section $R_Z:E_Z\to\B$ has
\[
 \M R_Zu(w)=\frac{g(w)r_Z(u)(w)}{h_Z(w)},\qquad
 qDR_Z=qR_ZM_w.\tag{T2}
\]
For completeness, the division operator giving (T2) is
$S_\rho F(x)=x^{-\rho}\int_x^\infty F(y)y^{\rho-1}dy$ when
$\M F(\rho)=0$. It satisfies $(D-\rho)S_\rho F=F$.
The zero moment rewrites its integral as minus the integral from zero
to $x$, proving rapid decrease at zero; the original formula proves it
at infinity. Iterating through the full multiplicities gives the
division by $h_Z$ in (T2). The source before division is
$\Theta(r_Z(u)(D)\phi_*)$, whose Mellin transform is $g r_Z(u)$.
Finally the remainder of multiplication of $r_Z(u)$ by $w$ is
$r_Z(M_wu)$ modulo $h_Z$, with quotient its original highest coefficient.
It follows that $DR_Z-R_ZM_w$ is a scalar multiple of $f_0$,
which lies in $\Theta V$. This proves the second assertion in (T2).

At zeros in $Z$, (T2) has exactly the jets $u$, by the unit $h_Z/g$;
at every other zero it vanishes to the full multiplicity. Hence
$J_{\rm all}R_Zu$ is the extension of $u$ by zero.
If $Z\subset Z'$ and $u'$ is extension by zero, then
\[
 r_{Z'}(u')=r_Z(u)\frac{h_{Z'}}{h_Z}.
\]
Indeed the right side has degree $<d_{Z'}$ and satisfies all defining
jets, including the new zero jets. Substitution in (T2) gives equal
Mellin transforms and thus equal reference sections (Mellin uniqueness
on any vertical line follows from Fourier uniqueness in $\log x$).
The $qR_Z$ therefore assemble into $\sigma$. Its full-jet left inverse
proves injectivity and $K\cap\sigma(E_{\rm fin})=0$.

An element of $Q$ maps to zero in $T$ exactly when its full-jet family
has finite support. Subtracting the corresponding $\sigma$ section
puts it in $K$. Thus the kernel before quotient is
$K+\sigma(E_{\rm fin})$, and after quotient it is canonically $K$.
This proves every assertion of exactness in (T1).

To see nonzero arithmetic tails directly, choose a nonzero nonnegative
$\chi\in C_c^\infty((1,2))$. Its Mellin transform $H(w)$ is entire
and is not identically zero, since $H(a)>0$ for real $a$.
For every actual zero $\rho$, the function $a\mapsto H(a+\rho)$
has only a discrete set of real zeros. The union of these exceptional
sets over the countable $Z_\infty$ is countable. Choose a real $a$
outside it and put $F(x)=x^a\chi(x)\in\B$.
Then every value coordinate of $J_{\rm all}F$ is nonzero. Since the
nontrivial zero set is infinite (the classical zero-counting theorem),
this family has infinite support and hence nonzero image in $T$.
\end{proof}

Here is the sheaf-level target explicitly. For $j:\C\hookrightarrow
\mathbb P^1(\C)$ and $\mathcal F(U)=\prod_{\rho\in U\cap Z_\infty}
E_\rho$, the old manuscript proves
\[
 0\to j_!\mathcal F\to j_*\mathcal F\to i_{\infty*}T\to0.\tag{T3}
\]
At a finite zero both stalks equal its full $E_\rho$; away from zeros
both vanish. At infinity the first stalk vanishes, and the second is
the product modulo finite modification: a neighborhood excludes only
finitely many points of the closed discrete support. This proves (T3)
on stalks. Globally $\Gamma(j_!\mathcal F)=E_{\rm fin}$ and
$\Gamma(j_*\mathcal F)=E_{\rm all}$, so (T1) is a receiver into the
actual global boundary of (T3). Neither $K=0$ nor $T_{\rm arith}=T$ is
asserted. This retains an analytic part and a global boundary that would
both be lost by identifying the programme with a finite zero divisor.

\section{Both endpoint states and the complete even-Schwartz source}
The earlier spectral-packet paper invokes the trace construction of Connes
and Consani. Their actual source range theorem, \cite[Lemma 4.1]{CCcycles},
retains both evaluation at zero and integration. Here is the full receiver
in the present programme's physical coordinate and with its factor two.
Let $\mathcal S_{\rm ev}=\mathcal S(\R)^{\rm ev}$, and choose a smooth
$\eta$ on $(0,\infty)$ equal to $1$ near zero and to $0$ for $x\geq2$.
Define the direct vector-space sum
\[
 \B_{01}=\B\oplus\C\eta\oplus\C x^{-1}\eta,
 \qquad b(F+c_0\eta+c_1x^{-1}\eta)=(c_0,c_1).
\]
Its definition is independent of $\eta$: two choices differ by an element
of $\B$, as do their products with $x^{-1}$.

\begin{theorem}[Exact two-endpoint receiver]\label{thm:two-endpoints}
The original theta sum extends to $\Theta:\mathcal S_{\rm ev}\to\B_{01}$,
and the following rows and boundary map are exact:
\[
 \begin{array}{ccccccccc}
 0&\to& V&\to&\mathcal S_{\rm ev}&\xrightarrow{a}&\C^2&\to&0\\
 &&\Theta\downarrow&&\Theta\downarrow&&\operatorname{diag}(-1,1)\downarrow\\
 0&\to&\B&\to&\B_{01}&\xrightarrow{b}&\C^2&\to&0,
 \end{array}
 \qquad a(\phi)=\left(\phi(0),\int_\R\phi\right).
 \tag{E1}
\]
Inclusion of the original target induces the canonical isomorphism
\[
 \B/\Theta V\ \xrightarrow{\ \sim\ }\
 \B_{01}/\Theta\mathcal S_{\rm ev}.\tag{E2}
\]
It intertwines the original $D=-x\partial_x$, and the action of $D$ on
each endpoint quotient $\C^2$ is $\operatorname{diag}(0,1)$.
The meromorphic Mellin transform of $F\in\B_{01}$ has precisely the
possible endpoint residues
\[
 \operatorname{Res}_{w=0}\M F=c_0,\qquad
 \operatorname{Res}_{w=1}\M F=c_1.\tag{E3}
\]
\end{theorem}
\begin{proof}
Poisson summation in the Fourier convention
$\widehat\phi(t)=\int_\R\phi(x)e^{-2\pi ixt}dx$ gives
\[
 \Theta\phi(x)=x^{-1}\!\int_\R\phi-\phi(0)
                         +x^{-1}\Theta\widehat\phi(1/x).
\]
The last term and all its $D$-derivatives decay more rapidly than every
power at zero. The defining theta series gives the corresponding decay
at infinity. This proves the middle arrow and the boundary matrix in
(E1). The source boundary map is onto: the Gaussian $g=e^{-\pi x^2}$
and $x^2g$ have boundary pairs $(1,1)$ and $(0,1/(2\pi))$.
The remaining exactness statements follow from the definitions.

To prove (E2) without discarding a boundary component, take any
$F\in\B_{01}$. Choose $\phi\in\mathcal S_{\rm ev}$ satisfying
$a(\phi)=\operatorname{diag}(-1,1)b(F)$. Then $F-\Theta\phi\in\B$,
proving surjectivity. If $F\in\B$ equals $\Theta\phi$ for some
$\phi\in\mathcal S_{\rm ev}$, the invertible boundary matrix forces
$a(\phi)=0$, so $\phi\in V$. This proves injectivity and proves
that a choice of boundary lift does not change the quotient map.

$D$ commutes with $\Theta$. Also $(D\phi)(0)=0$ and
$\int_\R D\phi=\int_\R\phi$ by integration by parts, while
$D1=0$ and $Dx^{-1}=x^{-1}$. Derivatives of the cutoff contribute
only elements of $\B$, proving the asserted quotient actions.
Finally subtract $c_0+c_1x^{-1}$ near zero inside the Mellin integral.
The remainder is entire in $w$, and integration over a sufficiently
small interval $(0,\epsilon)$ gives
$c_0\epsilon^w/w+c_1\epsilon^{w-1}/(w-1)$.
This proves (E3), retaining both signs.
\end{proof}

\begin{proposition}[Original Laplacian range, with explicit inverses]
The operator $D(D-1):\mathcal S_{\rm ev}\to V$ is a bijection.
For $\psi\in V$, its inverse is constructed on $x>0$ by
\[
 g(x)=\int_x^\infty\frac{\psi(t)}t\,dt,
 \qquad f(x)=x^{-1}\int_x^\infty g(t)\,dt,
 \tag{E4}
\]
followed by smooth even extension at zero. In particular
$\Theta D(D-1)e^{-\pi x^2}=f_0$, with exactly the original source
and its completed Mellin transform $2\xi$.
\end{proposition}
\begin{proof}
Since $\psi(0)=0$ and $\psi$ is even, $\psi(t)/t$ has a smooth odd
extension at zero. The first integral is smooth even there and rapidly
decreasing at infinity, and $Dg=\psi$. Fubini gives
$\int_0^\infty g(t)dt=\int_0^\infty\psi(t)dt=0$.
Consequently the second expression equals
$-x^{-1}\int_0^xg(t)dt$, which is smooth and even at zero.
Its expression as a tail integral proves Schwartz decay at infinity,
and $(D-1)f=g$. Thus $D(D-1)f=\psi$.
Conversely $D(D-1)f$ has value zero at zero and zero integral, so lies
in $V$. The equation $D(D-1)f=0$ on $(0,\infty)$ has general solution
$c+d/x$. Smoothness at zero and decay at infinity force $c=d=0$,
proving uniqueness.
\end{proof}

The precise unweighted/weighted trace dictionary is
$E\phi(x)=\tfrac12 x^{1/2}\Theta\phi(x)$ and
$H=x\partial_x=-D$; hence $H(1+H)=D(D-1)$ before conjugating
by multiplication by $x^{1/2}$. This is an operator identity, not an
identification of unrelated Hilbert metrics. The Connes--Consani spectral
quotient uses the closure of the source range in its specified
strong-Schwartz topology. For the original $\B$ topology there is always
the exact map
\[
 0\to\overline{\Theta V}/\Theta V\to Q
       \to\B/\overline{\Theta V}\to0.\tag{E5}
\]
Every original finite Mellin jet is continuous in that topology and
annihilates $\overline{\Theta V}$, so the finite sections (T2) survive
this passage. The full-source equality is now proved, rather than assumed:
$\overline{\Theta V}^{\mathscr B}=\Theta V$ in the topology
$p_{b,j}(F)=\sup_{x>0}x^b|D^jF(x)|$. Hence the first term of (E5)
is zero and the displayed quotient map is the identity isomorphism.
The complete proof, with the exact identity of seminorm systems and
continuous inverse estimates, is in the appended part
\emph{Closed range of the original full theta source}, or in
\texttt{ORIGINAL\_THETA\_CLOSED\_RANGE.tex}, Theorem
\texttt{thm:closed} and equation \texttt{eq:E5}; its human source is
Ralf Meyer, arXiv:math/0412277v3, Theorem 3.3. The two-leg source
generator is $D\oplus(1-D)$, with target generator $D$.
For a restricted source $W\subset V$ the exact replacement is
$\overline{\Theta W}^{\mathscr B}=\Theta\overline W^{V}$,
not an unproved density assertion. No Hilbert-space closedness is inferred.
For the Gaussian, (E3) gives residues $(-1,1)$ and the endpoint principal
part $1/(w(w-1))$. Thus the earlier quotient/lift residue calculation
and the current theta-source construction have the same exact two-endpoint
complex, with their common boundary explicitly removed by $D(D-1)$.

\section{The old local phase germ gives every entry of the original Taylor unit}
The atomic-phase manuscript retains, at a zero $\rho$ of exact multiplicity
$m_\rho$, the complete holomorphic factorization
\[
 \xi(\rho+z)=\lambda_\rho z^{m_\rho}
       \exp\!\left(\sum_{n\geq1}c_{\rho,n}z^n\right),\qquad
 \lambda_\rho=\frac{\xi^{(m_\rho)}(\rho)}{m_\rho!}.
 \tag{U0}
\]
Its phase-walk interpretation is explicitly credited there to unpublished
correspondence of Jacolm Tobley. The receiver here uses the full complex
logarithmic germ (U0); it makes no claim that the absolute packet weight
$|\lambda_\rho|$ determines this germ.

\begin{theorem}[All entries of the original unit, without removing phases]
Retain a finite actual packet divisor
$h_Z(w)=\prod_{\rho\in Z}(w-\rho)^{m_\rho}$ of degree $d$ and
$g=2\xi$. For $\rho\in Z$, define
\[
 H_\rho=\prod_{\sigma\in Z\setminus\{\rho\}}(\rho-\sigma)^{m_\sigma},
 \quad S_{\rho,n}=\sum_{\sigma\ne\rho}
                    \frac{m_\sigma}{(\rho-\sigma)^n},\quad
 a_{\rho,n}=\frac{(-1)^{n-1}}nS_{\rho,n}-c_{\rho,n}.
\]
The complete local jet of the original multiplier
$\varepsilon_Z=j_{h_Z}(h_Z/g)$ has coefficients
\[
 \varepsilon_{\rho,0}=\frac{H_\rho}{2\lambda_\rho},\qquad
 \varepsilon_{\rho,r}=\frac1r\sum_{n=1}^r
                  n a_{\rho,n}\varepsilon_{\rho,r-n}
       \quad(1\leq r<m_\rho).\tag{U1}
\]
Let $V_Z$ be the actual confluent evaluation matrix from the original
increasing power basis to the divided Taylor coefficients, namely
\[
 (V_Z)_{(\rho,j),\ell}=
 \begin{cases}\binom\ell j\rho^{\ell-j},&\ell\geq j,\\0,&\ell<j,
 \end{cases}
 \qquad 0\leq\ell<d.
\]
Let $T_\rho$ be the lower triangular matrix with
$(T_\rho)_{rj}=\varepsilon_{\rho,r-j}$ for $r\geq j$ and zero
otherwise. For the divided-jet input $u\in E_Z$ used in (T2), the exact
jet-to-power-coefficient map is
\[
 \operatorname{coeff}(r_Z(u))
   =V_Z^{-1}\!\left(\bigoplus_{\rho\in Z}T_\rho\right)u.
\]
If instead the input is the increasing-power coefficient vector $a$ of
the class whose jets are $u=V_Za$, the exact multiplication matrix is
\[
 \boxed{\quad M_{\varepsilon_Z}
     =V_Z^{-1}\!\left(\bigoplus_{\rho\in Z}T_\rho\right)V_Z.\quad}
 \tag{U2}
\]
Thus the two displayed matrices have different specified input coordinates;
the intervening morphism is precisely the confluent evaluation $V_Z$.
All root differences, full multiplicities and complex phases occur in
this original-coordinate formula. If $Z$ is stable under the two
functional-equation symmetries, then
\[
 \varepsilon_{1-\rho,r}=(-1)^{d+r}\varepsilon_{\rho,r},
 \qquad \varepsilon_{\bar\rho,r}=
                         \overline{\varepsilon_{\rho,r}}.\tag{U3}
\]
\end{theorem}
\begin{proof}
Every factor $\rho-\sigma$ is nonzero. On a sufficiently small disc,
the logarithm with value zero of
$\prod_{\sigma\ne\rho}(1+z/(\rho-\sigma))^{m_\sigma}$ has
coefficient $(-1)^{n-1}S_{\rho,n}/n$. Dividing this product by the
nonvanishing unit in (U0), and retaining the factor $2$ in $g=2\xi$,
gives
\[
 \frac{h_Z(\rho+z)}{g(\rho+z)}
   =\frac{H_\rho}{2\lambda_\rho}
                        \exp\!\left(\sum_{n\geq1}a_{\rho,n}z^n\right).
\]
Differentiate this convergent identity and compare coefficients of
$z^{r-1}$ to obtain (U1). Multiplication of truncated Taylor polynomials
is the Cauchy convolution of their coefficients, precisely $T_\rho$.
The map $V_Z$ is injective: a polynomial of degree less than $d$ with
all the stated zero jets is divisible by the monic degree-$d$ polynomial
$h_Z$ and hence vanishes. Equal finite dimensions give invertibility,
proving (U2). Finally, for a symmetry-stable divisor,
$h_Z(1-w)=(-1)^d h_Z(w)$, $g(1-w)=g(w)$, and both $h_Z,g$ commute
with conjugation. Apply these identities at $1-\rho+z$ and
$\bar\rho+z$ and compare Taylor coefficients to obtain (U3).
\end{proof}

This receiver converts a previously separate phase manuscript into actual
entries of the programme's canonical finite lift. In particular it is
available before any metric or determinant calculation: the lifted source
is exactly
$\Theta(r_Z(u)(D)\phi_*)$, divided by $h_Z(D)$ as in (T2).
The determinant of $M_{\varepsilon_Z}$ equals
$\prod_\rho\varepsilon_{\rho,0}^{m_\rho}$, but the full source and
its Gram cross terms retain the remaining coefficients in (U1).
No estimate on that determinant is substituted for an estimate on the
complete map (U2).

\section{The old jet-ideal chain enters Deligne's nilpotent filtration}
The earlier Boolean-jet and refined-divisor manuscripts prove that the
ideals of $\C[t]/t^m$ form the intrinsic chain
$I_j=(t^j)$, and give the adjunctions from this chain to the Boolean
multiplicity cube. Here this chain is carried into the original finite
theta class and its actual metric. The integer grading below is precisely
the nilpotent grading of Deligne \cite[(1.6.1),(1.6.7)]{DeligneII}; it
does not assign an arithmetic Frobenius weight to a zeta zero.

Retain $E_Z=\C[w]/h_Z$, its original increasing-power basis, and the
confluent matrix $V_Z$ of (U2). On its actual primary factor at $\rho$
write $t=w-\rho$ and $N_\rho u=tu$. Define
\[
 M_rE_\rho=
 \operatorname{span}\{t^j:0\leq j<m_\rho,\ m_\rho-1-2j\leq r\},
 \qquad M_rE_Z=V_Z^{-1}\!\bigoplus_\rho M_rE_\rho.
 \tag{N1}
\]
In particular the exponent of the ideal in the first formula is exactly
$\max(0,\min(m_\rho,\lceil(m_\rho-1-r)/2\rceil))$.
Let $N=V_Z^{-1}(\bigoplus N_\rho)V_Z$; this removes the semisimple
part by its actual primary idempotents, without changing $A=M_w$.

\begin{theorem}[Full multiplicity filtration and its original receiver]
The filtration (N1) satisfies
\[
 NM_r\subset M_{r-2},\qquad
 N^r:\operatorname{Gr}_r^M E_Z\xrightarrow{\sim}
             \operatorname{Gr}_{-r}^M E_Z\quad(r\geq0).
 \tag{N2}
\]
It is Deligne's canonical nilpotent filtration. It is preserved by the
original Taylor-unit multiplication $M_{\varepsilon_Z}$, whose action on
the grade belonging to $(\rho,j)$ is the full scalar
$\varepsilon_{\rho,0}=H_\rho/(2\lambda_\rho)$ from (U1).
The injective original section $\sigma:E_Z\to Q$ of (T2) carries it to
an $A$-stable filtration of $\sigma(E_Z)$. On that image the action is
the restriction of the original $D$, and $N$ is transported by the
actual spectral idempotents of that restriction.
\end{theorem}
\begin{proof}
Multiplication by $t$ lowers the displayed weight by two, proving the
inclusion. If $r\geq0$, the source grade at $\rho$ is zero unless
$r\leq m_\rho-1$ and $r\equiv m_\rho-1\pmod2$. In the remaining
case it is generated by $t^{(m_\rho-1-r)/2}$. Its image under $N^r$
is $t^{(m_\rho-1+r)/2}$, the nonzero generator of the target grade.
This proves every isomorphism with its coefficient equal to one in
the original divided Taylor coordinates. Deligne's uniqueness theorem
identifies this verified filtration with his canonical one.

For clarity, uniqueness on each actual cyclic primary can also be seen
directly. An $N_\rho$-stable subspace is an ideal of $\C[t]/t^m$,
so the weights $a_j$ at which $t^j$ enters any such filtration satisfy
$a_{j+1}\leq a_j-2$. The graded isomorphisms force the largest and
smallest weights to be $b$ and $-b$. The nonzero map $N^b$ from the
top grade implies $b\leq m-1$; the $m-1$ successive drops imply
$b\geq m-1$. Equality holds and every drop is exactly two, giving
$a_j=m-1-2j$. This proves uniqueness directly among the original
primary-stable filtrations as well.

In each primary, multiplication by the unit
$\sum\varepsilon_{\rho,j}t^j$ preserves every ideal $(t^j)$.
Its inverse is again a unit, so preservation is equality. Modulo the
next nonzero ideal step only its constant term survives. This proves
the asserted full scalar, including its phase and root-difference
factors. The identity $qDR_Z=qR_ZA$, proved in (T2), transports the
entire filtration into the original quotient. Since the different
primary eigenvalues are distinct, their idempotents are the actual
Chinese-remainder polynomials in $A$; applying those same polynomials
in $D|_{\sigma(E_Z)}$ gives exactly the stated transported $N$.
\end{proof}

There is an explicit raising operator on the associated graded, without
introducing a replacement metric:
\[
 R_\rho t^j=j(m_\rho-j)t^{j-1},\qquad
 H_\rho^{\rm gr}t^j=(m_\rho-1-2j)t^j.
 \tag{N3}
\]
The value is zero for $j=0$. Direct substitution gives
$[H^{\rm gr},N]=-2N$, $[H^{\rm gr},R]=2R$, and
$[R,N]=H^{\rm gr}$. For example the last diagonal coefficient is
$(j+1)(m-j-1)-j(m-j)=m-1-2j$, including the two endpoint cases.
Thus the old chain supplies the exact $\mathfrak{sl}_2$ grading used
in Deligne (1.6); these formulas do not assert $R=N^*$ for the original
arithmetic metric.

\begin{proposition}[Every graded quotient retains its original metric]
For this metric assertion let $h_{\rm act}$ be any of the programme's
actual monic finite root polynomials, retaining every root $\kappa$
and its full multiplicity $m_\kappa$. Write
$E_{\rm act}=\C[w]/h_{\rm act}$ and let $V_{\rm act}$ be its own
confluent evaluation matrix, exactly as in (U2). Formula (N1), with
these roots and multiplicities, defines its nilpotent filtration.
Let $G_N$ be its existing positive definite attained arithmetic metric
in the original increasing-power coefficients. Thus this assertion
applies directly to the full sum polynomial and $G_{q-1}$ of CMO1;
it does not identify the sum roots with zeros of $\xi$ or apply (U1)
to a polynomial that does not divide $2\xi$.
Order the pairs $(\kappa,j)$ by their actual weights
$m_\kappa-1-2j$, using a fixed order inside ties, and let $C$ be
the matrix with columns $V_{\rm act}^{-1}e_{\kappa,j}$ in this order.
Put $\mathcal H=C^*G_NC$. For a weight $r$ write $<$ for all earlier
grades and $r$ for the current grade. The metric on the actual quotient
$M_r/M_{r-1}$ is exactly
\[
 \mathcal G_r=\mathcal H_{rr}
      -\mathcal H_{r<}\mathcal H_{<<}^{-1}\mathcal H_{<r}.
 \tag{N4}
\]
For the first grade the subtracted term is zero. Moreover
\[
 \det G_N=|\det V_{\rm act}|^2\prod_r\det\mathcal G_r.
 \tag{N5}
\]
For the exact induced map $T_r=N^r:\operatorname{Gr}_r\to
\operatorname{Gr}_{-r}$, its squared norm on a coefficient vector $u$
is $(T_ru)^*\mathcal G_{-r}(T_ru)$, and its domain squared norm is
$u^*\mathcal G_ru$. Thus its actual squared singular values are the
eigenvalues of
\[
 \mathcal G_r^{-1/2}T_r^*\mathcal G_{-r}T_r
                  \mathcal G_r^{-1/2}.
 \tag{N6}
\]
\end{proposition}
\begin{proof}
For fixed current coefficient $u$, all lifts in $M_r$ are $(v,u)$
with lower-grade coefficient $v$. Their squared norms are
\[
 v^*\mathcal H_{<<}v+v^*\mathcal H_{<r}u+
 u^*\mathcal H_{r<}v+u^*\mathcal H_{rr}u.
\]
Completing this exact quadratic form gives the unique minimum at
$v=-\mathcal H_{<<}^{-1}\mathcal H_{<r}u$ and the value
$u^*\mathcal G_ru$, proving (N4). The positive definiteness of
$\mathcal H$ proves the stated inverse exists and $\mathcal G_r>0$.
Block elimination with these same coefficients yields
$\det\mathcal H=\prod_r\det\mathcal G_r$ inductively. Since
$C$ is $V_{\rm act}^{-1}$ followed by a permutation,
$\det\mathcal H=|\det V_{\rm act}|^{-2}\det G_N$, proving (N5) with the
full original root-difference determinant retained. The induced map
$T_r$ is the one already computed in (N2); writing its two quadratic
forms proves (N6). No cross term was replaced by a rank count.
\end{proof}

This gives a fully specified recipient for Deligne's nilpotent
filtration within the old split-zero jet geometry. It also identifies
the exact analytic matrices that a mixed-state estimate has to control.
The integers in (N1) are nilpotent grades; arithmetic purity would
add an assertion about eigenvalue magnitudes, which is neither a
consequence of these integers nor included in this result.

\section{Repair of the earlier finite-prime boundary lift}
The earlier primitive-defect and split-zero-defect manuscripts use the
finite-prime extension below in their semilocal arithmetic boundary
construction. Their proposed invariant scalar cutoff is not by itself
an element of the crossed product: the quotient unit is a rank-one
Morita projection, rather than a constant function in $C_0(\R)$.
The following supplies that projection, its actual self-adjoint lift,
and its entire connecting-map calculation. The sign convention remains
the source's increasing logarithmic coordinate.

Fix a prime $p$, put $\ell=\log p$, and retain
\[
 Y=(\mathbb Z\cup\{\infty\})\times\R,\qquad
 T(m,y)=(m+1,y+\ell),\quad T(\infty,y)=(\infty,y+\ell),
\]
where $m\to+\infty$ converges to the adjoined point. Set
\[
 I=C_0(\mathbb Z\times\R)\rtimes_T\mathbb Z,
 \quad E=C_0(Y)\rtimes_T\mathbb Z,
 \quad Q_p=C_0(\R)\rtimes_{y\mapsto y+\ell}\mathbb Z.
 \tag{K1}
\]
Thus $0\to I\to E\to Q_p\to0$ is the original extension. Use the
crossed-product convention $U f U^{-1}=f\circ T^{-1}$, so that for
finite coefficient families
\[
 (a*b)_n(z)=\sum_k a_k(z)b_{n-k}(T^{-k}z),\qquad
 (a^*)_n(z)=\overline{a_{-n}(T^{-n}z)}.
 \tag{K2}
\]

Choose a nonnegative $c\in C_c(\R)$ positive on $[-\ell,\ell]$ and
put
\[
 \chi(y)=\frac{c(y)}{\sqrt{\sum_{n\in\mathbb Z}c(y-n\ell)^2}}.
\]
The denominator is a strictly positive continuous periodic function;
all sums are locally finite. Hence $\chi\in C_c(\R)$ and
$\sum_n\chi(y-n\ell)^2=1$ at every $y$. Define
\[
 P_n(y)=\chi(y)\chi(y-n\ell).
 \tag{K3}
\]
Only finitely many $n$ occur, and each coefficient has compact support.
Let $\phi(t)=(1+e^t)^{-1}$ and define
\[
 H_n(m,y)=\phi(y-m\ell)P_n(y),\qquad
 H_n(\infty,y)=P_n(y).
 \tag{K4}
\]

\begin{theorem}[Actual finite-prime connecting morphism]
The family $P$ is a full projection in $Q_p$ whose corner is
$P Q_pP\cong C(\R/\ell\mathbb Z)$; it represents the unit class
under the original circle Morita equivalence. The family $H$ is an
element of $E$, is self-adjoint with $0\leq H\leq1$, and lifts $P$.
Under the explicit identification $I\cong C_0(\R,\mathcal K(\ell^2\mathbb Z))$,
the exponential $\exp(2\pi iH)$ has Fredholm determinant
\[
 \det\bigl(\exp(2\pi iH)(t)\bigr)=e^{2\pi i\phi(t)}.
\]
Consequently the original connecting homomorphism, with
counterclockwise winding positive, is
\[
 \boxed{\partial_p([P])=-1.}\tag{K5}
\]
\end{theorem}
\begin{proof}
The involution formula (K2) gives $P^*=P$, and
\[
 (P*P)_n(y)=\chi(y)\chi(y-n\ell)
                 \sum_k\chi(y-k\ell)^2=P_n(y),
\]
so $P$ is a projection in the actual quotient algebra. The regular
representation over an orbit $y+\ell\mathbb Z$ sends it to the
rank-one projection onto the unit vector
$v_y=(\chi(y+r\ell))_{r\in\mathbb Z}$. If $y$ is increased by
$\ell$ these vectors and the representations are related by the
same bilateral shift. Compression by $v_y$ therefore gives a scalar
continuous function on $\R/\ell\mathbb Z$. Conversely every continuous
periodic scalar function times $P$ belongs to the corner, giving the
claimed isomorphism. The projection is full: locally on this compact
orbit space its nonzero continuous unit vector generates the rank-one
operators, and finite-rank matrices generated from these are dense in
the compact-operator fibers. A finite partition of unity on the circle
then gives the whole algebra as the closed ideal generated by $P$.
This is the concrete free-proper-action Morita construction of Green
\cite{Green}, not an identification of a constant multiplier with an
algebra projection.

For each nonzero coefficient in (K4), $y$ stays in the fixed compact
support of $\chi$. As $m\to+\infty$, $y-m\ell\to-\infty$ uniformly
on that support, and $\phi(y-m\ell)\to1$, proving continuity at
the adjoined stratum. As $m\to-\infty$ it tends uniformly to zero.
It follows that $H_n\in C_0(Y)$. The integer $n$-support is finite,
so $H\in E$. The scalar $t=y-m\ell$ is invariant under $T$ on
the generic part. Consequently (K2) proves self-adjointness and
\[
 (H^a)_n(m,y)=\phi(t)^a P_n(y)\qquad(a\geq1),
\]
with the corresponding boundary value $P_n(y)$.

For completeness, the generic algebra identification is given on the
dense finite convolution algebra by
\[
 (\pi_t(a))_{r,s}=a_{r-s}(r,t+r\ell).
 \tag{K6}
\]
Substituting (K2) verifies multiplication and adjoints entry by entry.
If a coefficient has finite $m$-support and compact $y$-support,
its image is a compactly supported continuous finite matrix family.
Conversely a family $f(t)e_{r,s}$ is obtained from the single coefficient
$a_{r-s}(m,y)=\mathbf1_{\{r\}}(m)f(y-r\ell)$. These elements are dense
on both sides. Thus the usual crossed-product completion gives the
displayed isomorphism (equivalently the explicit free-proper-action
case of \cite{Green}). Under it
\[
 \pi_t(H)=\phi(t)\,v_tv_t^*,\qquad \|v_t\|=1.
\]
At every boundary orbit the representation gives the projection $P$.
These fiber representations detect positivity; hence $0\leq H\leq1$.

The exponential formula is now exact, with no missing cross terms:
\[
 \pi_t(\exp(2\pi iH))
   =I+(e^{2\pi i\phi(t)}-1)v_tv_t^*.
 \tag{K7}
\]
The scalar coefficient tends to zero at both ends of $\R$, so this
is a unitary in the unitization of the original ideal. Its rank-one
Fredholm determinant is the asserted scalar function. That determinant
represents its $K_1(C_0(\R)\otimes\mathcal K)$ class: one can deform
$v_t$ to a fixed unit vector by normalized straight-line interpolation,
since $v_t$ and any fixed coordinate unit vector are nonnegative and
the interpolation never vanishes. The coefficient of the rank-one
projection vanishes at the two ends uniformly throughout this homotopy.
This supplies an explicit homotopy to a fixed rank-one loop with the
same determinant.

By the exponential definition of the boundary map, $\partial_p[P]$
is represented by this unitary. When $t$ increases from $-\infty$ to
$+\infty$, $\phi$ decreases from one to zero. The determinant therefore
makes exactly one clockwise circuit. Its counterclockwise degree is
$-1$, proving (K5).
\end{proof}

\subsection*{The actual visible semilocal receiving map}
The earlier source specifies the visible semilocal space after quotient
by the finite local unit groups. Its map to (K1) can also be given
explicitly. Let $S=\{p_1,\ldots,p_n,\infty\}$, put
$\ell_i=\log p_i$, and let
\[
 X_S=\{(m_1,\ldots,m_n,x)\in
    (\mathbb Z\cup\{\infty\})^n\times\R:
    \text{at most one of }m_i=\infty\text{ or }x=0\}.
\]
It carries the subspace topology from the displayed product. The group
$\Gamma=\mathbb Z^n\times\{\pm1\}$ acts by
\[
 (a,\epsilon)(m,x)=
       (m+a,\epsilon e^{\sum_i a_i\ell_i}x),\qquad
 \infty+a_i=\infty.
 \tag{K8}
\]
This is precisely the action of rational $S$-units on the local
valuations and real coordinate after finite local unit quotient. Let
$A_S=C_0(X_S)\rtimes\Gamma$, let $A_\varnothing$ be the ideal of
its all-nonzero stratum, and let $B_S$ be the quotient algebra of the
one-zero strata.

\begin{theorem}[Full primitive boundary row for the specified space]
The extension $0\to A_\varnothing\to A_S\to B_S\to0$ has, with
the original logarithmic orientation,
\[
 \delta_0:\mathbb Z^{n+2}\longrightarrow\mathbb Z,
 \qquad
 \delta_0(k_1,\ldots,k_n,k_+,k_-)
       =-\sum_i k_i-k_+-k_-.
 \tag{K9}
\]
In particular
$K_0(A_S)\cong\ker\delta_0\cong\mathbb Z^{n+1}$ and
$K_1(A_S)\cong\mathbb Z^n$. Each finite-prime component is received
from (K5) by an explicit equivariant coordinate map, and the two
archimedean components each have the same sign $-1$.
\end{theorem}
\begin{proof}
On the invariant open edge containing the generic stratum and the
$p_i$-zero stratum, all $m_j$ for $j\ne i$ are finite and $x\ne0$.
Use the actual coordinates
\[
 z_j=m_j\ (j\ne i),\quad m=m_i,\quad
 y=\log|x|-\sum_{j\ne i}m_j\ell_j,\quad s=\operatorname{sign}x.
 \tag{K10}
\]
They give a homeomorphism to
$\mathbb Z^{n-1}\times Y_{p_i}\times\{\pm1\}$.
The inverse is
$m_j=z_j$ for $j\ne i$, $m_i=m$, and
$x=s\exp(y+\sum_{j\ne i}z_j\ell_j)$, proving the assertion also
at $m_i=\infty$. Under (K8), the other prime generators translate
only their $z_j$, the sign generator translates only $s$, and the
$p_i$ generator acts only by $(m,y)\mapsto(m+1,y+\ell_i)$.
Taking crossed products therefore gives the literal tensor factor
\[
 \mathcal K(\ell^2\mathbb Z^{n-1})\otimes E_{p_i}\otimes M_2(\C),
\]
with matching generic ideal and boundary quotient. The compact factor
identification follows directly from the matrix units for translation
of a discrete group on itself. Removing these compact factors is the
explicit Morita map, and the original generic logarithmic coordinate is
\[
 t=\log|x|-\sum_jm_j\ell_j=y-m_i\ell_i.
\]
Thus the orientation is exactly the one in (K5). Naturality of the
exponential boundary construction yields the component $-1$.

On the invariant archimedean edge every $m_i$ is finite and $x$ may
be zero. Use instead
$u=x\exp(-\sum_i m_i\ell_i)$. This is an equivariant homeomorphism
to $\mathbb Z^n\times\R$: prime generators translate the first
factor only, and the sign acts by $u\mapsto-u$. Hence its extension
is a compact-operator tensor factor times
\[
 0\to C_0(\R^\times)\rtimes\{\pm1\}
 \to C_0(\R)\rtimes\{\pm1\}
 \to C^*(\{\pm1\})\to0.
\]
Write $p_\pm=(1\pm U)/2$ in the quotient. Choose an even continuous
real-valued $\psi(u)=e^{-u^2}$, with $\psi(0)=1$ and
$\psi(u)\to0$ as $|u|\to\infty$.
The actual self-adjoint lifts are $\psi p_\pm$ in this crossed
product. On $u>0$ the orbit representation gives the rank-one matrices
$P_\pm=\frac12\left(\begin{smallmatrix}1&\pm1\\\pm1&1\end{smallmatrix}\right)$.
The exponential determinant is $e^{2\pi i\psi(u)}$ for either
choice. With $t=\log u$, it goes from phase $2\pi$ to zero as $t$
goes from $-\infty$ to $+\infty$, giving $-1$ for both classes.
This logarithmic coordinate is again exactly the common generic
coordinate above.

The boundary consists of the disjoint $n$ finite-prime strata and
the archimedean stratum. It has groups
$K_0(B_S)=\mathbb Z^{n+2}$ and $K_1(B_S)=\mathbb Z^n$ by the
circle corners and the two sign projections just constructed.
The generic ideal is stably $C_0(\R)$, so its $K_0$ is zero and
its $K_1$ is $\mathbb Z$. The extension and all edge inclusions have
the same generic ideal; restriction to each summand of $K_0(B_S)$
therefore gives exactly the computed component. This proves (K9).
Since the row is surjective, the six-term exact sequence gives
$K_0(A_S)=\ker\delta_0$ and $K_1(A_S)\cong K_1(B_S)$, proving the
remaining assertions. For $n=1$ the kernel is the actual primitive
$A_2$ lattice in $\mathbb Z^3$, with its three oriented generators.
\end{proof}

Thus the actual model calculation and the visible semilocal receiving
maps are both complete. This is a theorem for the specified original
finite-unit-quotiented space $X_S$, not a claim about an unspecified
crossed product. It restores the complete boundary morphism on this
earlier branch of the Split Zero RH programme, while leaving arithmetic
metric positivity and eigenvalue bounds to their genuine analytic
calculations.

\section{Denominator-history inversion reconstructs the original source}
The earlier arithmetic-curve manuscript's incidence-inversion theorem
is the exact finite divisor identity
$\sum_{d\mid n}\mu(d)=\mathbf1_{\{n=1\}}$. Its analytic receiver
in the original physical coordinate supplies a useful description of
the entire source quotient, including the part not detected by a fixed
finite packet. No coordinate or source measure is changed. In this
section the same theta formula is also defined on the space of smooth
functions on $(0,\infty)$ that are rapidly decreasing with all
$D$-derivatives at infinity and are $O(x^{-1})$ with those derivatives
at zero. Its defining series converges locally on this space; its
restriction to the original $V$ is the original $\Theta$.

\begin{theorem}[Exact original-source reconstruction]
For $F\in\B$ define
\[
 (\mathcal I_\mu F)(x)=\frac12\sum_{n\geq1}\mu(n)F(nx)
       \qquad(x>0).\tag{MI1}
\]
The sum and all its $D$-derivatives converge absolutely locally
uniformly on $(0,\infty)$. Its value is smooth, rapidly decreasing
with all $D$-derivatives at infinity, and is $O(x^{-1})$ with all
$D$-derivatives at zero. It satisfies
\[
 \Theta\mathcal I_\mu F=F,\qquad
 \mathcal I_\mu\Theta\phi=\phi\quad(\phi\in V),\qquad
 D\mathcal I_\mu=\mathcal I_\mu D.
 \tag{MI2}
\]
Consequently the original algebraic source quotient has the exact
operator-compatible presentation
\[
 \boxed{\quad Q=\B/\Theta V\ \cong\ \mathcal I_\mu(\B)/V.\quad}
 \tag{MI3}
\]
The obstruction on the right is membership in the original even
Schwartz zero-moment space, not inversion of the theta sum on positive
$x$. On every line $\Re w>1$,
\[
 \mathcal M(\mathcal I_\mu F)(w)
       =\frac{\mathcal MF(w)}{2\zeta(w)}.
 \tag{MI4}
\]
\end{theorem}
\begin{proof}
For each compact positive $x$-interval and every $r$, the summand
$D^rF(nx)$ is bounded by $C_B n^{-B}$ for every $B>1$. This proves
absolute local convergence and differentiation. For $x\geq1$ the
same bound is $C_B n^{-B}x^{-B}$, proving rapid decrease.
For $g=D^rF$, both $g$ and its ordinary derivative are integrable
on $(0,\infty)$. On the interval $((n-1)x,nx)$ the fundamental
theorem of calculus gives
\[
 |g(nx)|\leq x^{-1}\int_{(n-1)x}^{nx}|g(y)|\,dy
                 +\int_{(n-1)x}^{nx}|g'(y)|\,dy.
\]
Sum over $n$ to obtain
$|D^r\mathcal I_\mu F(x)|\leq
 (\|g\|_1/x+\|g'\|_1)/2$, proving the zero-end estimate.

For fixed $x>0$, the double sum involved in either composition is
absolutely convergent: its absolute value is bounded by
$\sum_{k\geq1}d(k)|F(kx)|$ or the analogous sum for $\phi$.
Use $d(k)\leq k$ and rapid decrease to verify convergence.
Regrouping by $k=mn$ then gives (MI2) by the classical finite divisor
identity $\sum_{d\mid k}\mu(d)=\mathbf1_{k=1}$, documented in
\cite[Chapter 2]{ApostolMobius} and NIST DLMF (27.5.2).
The derivative statement follows from the already proved locally
uniform convergence. In particular $V\subset\mathcal I_\mu(\B)$,
and (MI2) makes the two induced maps in (MI3) inverse. This proves
both injectivity and surjectivity without replacing the original $V$.
Finally absolute integrability on $\Re w>1$ permits substitution
$y=nx$ in each Mellin integral. The absolutely convergent Euler
identity $\sum_n\mu(n)n^{-w}=1/\zeta(w)$ proves (MI4).
\end{proof}

\begin{proposition}[Complete primary source singularities of the actual lift]
For the original $R_Zu$ of (T2), put
$\psi_u=\mathcal I_\mu R_Zu$. Its exact equation and Mellin transform
are
\[
 h_Z(D)\psi_u=r_Z(u)(D)\phi_*,\qquad
 \mathcal M\psi_u(w)=
       \frac{r_Z(u)(w)}{h_Z(w)}\,\mathcal M\phi_*(w),
 \tag{MI5}
\]
where
$\mathcal M\phi_*(w)=\frac12w(w-1)\pi^{-w/2}\Gamma(w/2)$.
Write the proper rational function in its unique original partial
fractions
\[
 \frac{r_Z(u)(w)}{h_Z(w)}
 =\sum_{\rho\in Z}\sum_{j=1}^{m_\rho}
                     \frac{b_{\rho,j}}{(w-\rho)^j}.
\]
Then, without discarding any primary or multiplicity,
\[
 \psi_u(x)=\sum_{\rho,j}\frac{b_{\rho,j}x^{-\rho}}{(j-1)!}
  \int_x^\infty \phi_*(y)y^{\rho-1}
                    \bigl(\log(y/x)\bigr)^{j-1}\,dy.
 \tag{MI6}
\]
Its full singular part at zero is
\[
 \sum_{\rho,j}\frac{b_{\rho,j}x^{-\rho}}{(j-1)!}
 \sum_{a=0}^{j-1}\binom{j-1}{a}
        (-\log x)^{j-1-a}(\mathcal M\phi_*)^{(a)}(\rho).
 \tag{MI7}
\]
The exact remainder obtained by subtracting (MI7) is $O(x^2)$ at
zero with all $D$-derivatives and has a convergent even Taylor series
there. Also $\int_0^\infty\psi_u(x)\,dx=0$. Thus (MI6) retains
the original zero-integral condition while explicitly displaying its
failure of smoothness at zero for every nonzero finite class.
\end{proposition}
\begin{proof}
Apply (MI2) and (MI4) to (T2). The identity
$g/(2\zeta)=\mathcal M\phi_*$ gives (MI5), including the factor
$1/2$. Iterating the integral operator $S_\rho$ in the proof of
(T2), now without a vanishing Mellin moment at $\rho$, gives
\[
 S_\rho^j\phi_*(x)=\frac{x^{-\rho}}{(j-1)!}
       \int_x^\infty\phi_*(y)y^{\rho-1}
                      (\log(y/x))^{j-1}\,dy.
\]
This follows by Fubini over $x<y_1<\cdots<y_j$: the ordered
simplex in the logarithmic variables has volume
$(\log(y_j/x))^{j-1}/(j-1)!$. Its Mellin transform on
$\Re w>1$ is $\mathcal M\phi_*(w)/(w-\rho)^j$.
The proper partial fractions therefore give (MI6) by Mellin
uniqueness on that line.

The integral from zero to infinity converges for each actual
$0<\Re\rho<1$, since $\phi_*(y)=O(y^2)$ at zero and decreases
rapidly at infinity. Expanding its logarithmic power gives exactly
(MI7). The remaining integral from zero to $x$ is bounded by a
constant times $x^{\Re\rho+2}$ before multiplication by
$x^{-\rho}$; substituting $y=xv$ proves the same bound for every
$D$-derivative. More explicitly, if
$\phi_*(x)=\sum_{n\geq1}c_{2n}x^{2n}$ near zero, that remainder
for $S_\rho^j\phi_*$ is precisely
\[
 (-1)^j\sum_{n\geq1}\frac{c_{2n}}{(\rho+2n)^j}x^{2n},
\]
convergent near zero. The signs follow from integrating
$(\log v)^{j-1}$ on $0<v<1$.

Formula (MI6) is integrable at zero because every actual
$\Re\rho<1$, including all logarithmic powers, and at infinity
by the tail estimates. Evaluation of its Mellin transform at $w=1$
is legitimate. Since $h_Z(1)\ne0$ and
$\mathcal M\phi_*(1)=0$, its integral is zero. Finally if
$\psi_u$ were in $V$, (MI2) would imply $R_Zu\in\Theta V$,
and its full original jets force $u=0$ by (T2). Thus no nonzero
primary class has been erased by this reconstruction.
\end{proof}

\begin{corollary}[Exact original $dx$ integrability threshold]
The reconstructed finite source $\psi_u=\mathcal I_\mu R_Zu$
satisfies
\[
 \boxed{\quad\psi_u\in L^2(\R_{>0},dx)
 \quad\Longleftrightarrow\quad
 u_\rho=0\ \text{for every }\rho\in Z\text{ with }
                         \Re\rho\geq\tfrac12.\quad}\tag{MI9}
\]
Here $u_\rho=0$ means the entire primary jet, not merely its value.
\end{corollary}
\begin{proof}
The local factor $\mathcal M\phi_*(\rho)$ is nonzero at every
nontrivial zero: its explicit Gamma formula has no zero or pole
there. Thus a principal part of
$\mathcal M\psi_u=(\mathcal MR_Zu)/(2\zeta)$ at $\rho$ is zero
exactly when the whole jet $u_\rho$ vanishes. Formula (MI7) shows
that each remaining primary contributes a nonzero polynomial in
$-\log x$ times $x^{-\rho}$, and the remainder is $O(x^2)$.
If all their real parts are less than $1/2$, every term is square
integrable at zero. Rapid decay at infinity proves sufficiency.

For necessity put $t=-\log x$. The singular part after multiplication
by the square-root density $e^{-t/2}$ is the finite exponential
polynomial $\sum_\rho e^{(\rho-1/2)t}P_\rho(t)$.
Let $b$ be the largest real part among its nonzero terms and $d$
the largest polynomial degree at that real part. After division by
$e^{(b-1/2)t}t^d$, the leading part is
$A(t)=\sum_{\Re\rho=b,\deg P_\rho=d}a_\rho e^{i\Im\rho t}$,
and the remaining terms tend uniformly to zero on every interval
$[T,T+L]$ as $T\to\infty$ for fixed $L$. At least one $a_\rho$
is nonzero. The distinct frequencies imply
\[
 \int_T^{T+L}|A(t)|^2dt
 =L\sum|a_\rho|^2+E(T,L),\qquad
 |E(T,L)|\leq
  \sum_{\rho\ne\sigma}\frac{2|a_\rho a_\sigma|}
                                    {|\Im\rho-\Im\sigma|}.
\]
Choose $L$ so the right-hand error is at most half the positive
diagonal term. This lower bound is uniform in $T$. The vanishing
remainder preserves a fixed positive fraction of it for all large
$T$. If $b\geq1/2$, multiplication back by
$e^{(b-1/2)t}t^d$ yields a positive lower bound on each of infinitely
many disjoint such intervals, and indeed an increasing bound unless
$b=1/2,d=0$. Their squared integrals diverge in every case. The
$O(x^2)$ remainder is square integrable and cannot change that
conclusion. This proves necessity without treating different
frequencies as pointwise orthogonal.
\end{proof}

There is a useful bounded tail estimate on its stated domain only.
For $\sigma>1$ in the weighted norm
$\|F\|_\sigma^2=\int_0^\infty|F(x)|^2x^{2\sigma-1}dx$,
dilation has norm $n^{-\sigma}$. The triangle inequality gives
\[
 \left\|\mathcal I_\mu F-\tfrac12\sum_{n\leq L}\mu(n)F(n\,\cdot)
 \right\|_\sigma
 \leq\frac{L^{1-\sigma}}{2(\sigma-1)}\|F\|_\sigma.
 \tag{MI8}
\]
The original $dx$ metric has $\sigma=1/2$; (MI8) is not a bound
in that metric. The exact maps (MI1)--(MI7), including both endpoint
regions, do not require such an unsupported substitution.

\section{The original theta complex over the earlier split-support space}
The earlier generic-point-suspension theorem identifies the underlying
space $X=\operatorname{Spec}G(\C)$ with two points:
\[
 \eta=\{\tau\},\qquad z=\{\tau,0_\C\},\qquad
 \operatorname{Open}(X)=\{\varnothing,\{\eta\},X\}.
\]
The new generic point is $\eta$ and the supported closed point is $z$.
Work here with sheaves of complex vector spaces on this underlying
space; the coefficient category is specified explicitly. In particular
no map making its Boolean local semiring into a complex algebra is
assumed. The original infinite-dimensional spaces give algebraic
sheaves, while the finite primary recipients below are constructible
finite-dimensional sheaves.

\begin{theorem}[The actual theta quotient is supported cohomology]
Define $\mathcal F_\Theta$ by
\[
 \mathcal F_\Theta(X)=V,\qquad
 \mathcal F_\Theta(\{\eta\})=\B,\qquad
 \operatorname{res}_{X,\{\eta\}}=\Theta.
 \tag{S1}
\]
Then its local cohomology with support at $z$ is calculated by the
original complex, with its original coefficient spaces and operator:
\[
 R\Gamma_{\{z\}}(X,\mathcal F_\Theta)\simeq[V\xrightarrow{\Theta}\B]
       \quad\text{in degrees }0,1.
 \tag{S2}
\]
Thus $H^0_{\{z\}}=0$, $H^1_{\{z\}}=Q$, and all higher groups vanish.
For every finite actual zero packet $Z$, its full Mellin-jet map and
canonical section give sheaf maps
\[
 \mathcal F_\Theta\xrightarrow{(0,J_Z)} j_!E_Z
           \xrightarrow{(0,R_Z)}\mathcal F_\Theta,
 \tag{S3}
\]
whose induced cohomology maps are exactly
$\overline J_Z:Q\to E_Z$ and $\sigma_Z:E_Z\hookrightarrow Q$.
Their composite on $E_Z$ is the identity, and both cohomology maps
intertwine the original $D$ and multiplication by $w$.
\end{theorem}
\begin{proof}
Every open cover of $X$ contains $X$ itself, so a sheaf of vector
spaces on $X$ is precisely an arrow $F_z\to F_\eta$; morphisms
are commutative squares. Its sections supported at $z$ are the
kernel of that arrow. Evaluation at either point is exact. The
objects $(A\to0)$ and $(B\xrightarrow{1}B)$ are injective: morphisms
into them identify respectively with $\operatorname{Hom}(F_z,A)$
and $\operatorname{Hom}(F_\eta,B)$, exact functors of vector spaces.
For a general arrow $d:A\to B$, the explicit injective resolution is
\[
 0\to(A\xrightarrow d B)
 \xrightarrow{(a\mapsto(a,da),\,1)}
 (A\oplus B\xrightarrow{\mathrm{pr}_B}B)
 \xrightarrow{((a,b)\mapsto da-b,\,0)}(B\to0)\to0.
 \tag{S4}
\]
Each component is an exact sequence, proving exactness of the sheaf
sequence. The middle term is the direct sum of the two kinds of
injective just identified. Applying the support-section functor
gives exactly $[A\xrightarrow d B]$. With $d=\Theta$, this proves
(S2). Its degree-zero group vanishes by the actual inverse in (MI2).
The same argument for $j_!E_Z=(0\to E_Z)$ gives its supported
cohomology $E_Z$ in degree one and zero elsewhere.

The first square in (S3) commutes because $J_Z\Theta=0$ with all
original multiplicities. The second square commutes because its
closed component has zero source. Their induced maps in (S4) are
the displayed original complex maps, proving the cohomology
identifications and $\overline J_Z\sigma_Z=1$. The first map
intertwines $D$ already on the complex. For the section, (T2) proves
$DR_Z-R_ZA=\Theta\gamma_Z$, where $\gamma_Zu$ is the original
scalar highest-coefficient term times $\phi_*$. Thus the failure
of strict intertwining is the explicitly supplied chain homotopy
through the original degree-zero source $V$; it vanishes on degree-one
cohomology. This proves the asserted operator equivariance without
changing the section.
\end{proof}

\begin{proposition}[The exact split semimodule underlying the same arrow]
Let $S=G(\C)$ with additive zero $\tau$ and supported scalar zero
$e=0_\C$. On the disjoint union
$M_\Theta=V\sqcup\B$, retain the ordinary vector addition within
each component, and define the mixed sum by
\[
 v\oplus b=b\oplus v=\Theta v+b\quad(v\in V,b\in\B).
\]
The zero of this monoid is $0_V$. A supported scalar $a\in\C$
acts by ordinary scalar multiplication within each component, and
$\tau$ sends every element to $0_V$. These operations make
$M_\Theta$ an $S$-semimodule. Its support semilattice is
\[
 eM_\Theta=\{0_V,0_\B\},\qquad0_V<0_\B,
\]
with transport map exactly $\Theta$. The map
\[
 \pi:M_\Theta\longrightarrow Q\sqcup\{\tau\},\qquad
 \pi(v)=\tau,\quad\pi(b)=[b],
 \tag{S5}
\]
is an $S$-linear surjection and identifies its target with the Bourne
quotient by the subsemimodule $V$. In particular the unsupported
zero fiber is $V$, whereas the supported-zero fiber is the separate
subset $\Theta V\subset\B$.
\end{proposition}
\begin{proof}
The support tag of a sum is the maximum of its two tags. Once a
$\B$ tag occurs, the value of any finite sum is the sum in $\B$
of its $\B$ terms and of the images under $\Theta$ of its $V$
terms. This formula proves associativity, commutativity and the
stated identity. Linearity of $\Theta$ proves distributivity of
supported scalars across every combination of tags. For a scalar
sum involving $\tau$, the scalar acts as the unchanged supported
scalar and the other summand contributes $0_V$, again the identity.
For a scalar product involving $\tau$, both sides are $0_V$.
The remaining scalar laws are the vector-space laws, including
the case of supported scalars whose sum is the supported zero $e$.
This verifies all semimodule axioms. Its two supported zero vectors
are distinct by the disjoint-union construction, and
$0_V\oplus0_\B=0_\B$.

Give $Q\sqcup\{\tau\}$ the usual split semimodule structure:
vector addition on $Q$, additive identity $\tau$, supported scalar
action on $Q$, and $\tau$ acting by the zero map to $\tau$.
Since $q\Theta=0$, (S5) preserves the mixed sum; it preserves all
other operations directly. Surjectivity follows from that of $q$.
In the Bourne relation, $x\sim y$ means
$x\oplus v=y\oplus v'$ for some $v,v'\in V$.
Both sides have the same support tag only when $x,y$ have the same
tag. All elements of $V$ become equivalent by their additive
inverses. For $b,b'\in\B$ the relation is precisely
$b-b'\in\Theta V$. Thus its equivalence classes are one unsupported
class and the supported vector quotient $Q$, proving the exact
quotient claim and both fibers. The original $D$ on both components
preserves these operations because $D\Theta=\Theta D$, so it also
descends to the original operator on $Q$.
\end{proof}

The earlier support-diagram theorem therefore has a concrete receiver
on the entire original theta complex, and the earlier split-spectrum
theorem has a concrete supported-cohomology receiver on the same
arrow. These are explicit functors into different stated coefficient
categories. The sheaf (S1) is not silently obtained by quasi-coherent
semiring localization: its generic complex vector space and the
Boolean stalk of the structure semiring have different types. The
finite maps (S3), the chain homotopy, and the split quotient (S5)
specify exactly how the programme's analytic and support data meet.

\section{Actual reflection, primary signs, and the unweighted Hilbert completion}
The compactified-shadow manuscript supplies the full-multiplicity
divisor and its Klein-four symmetry. We now realize that symmetry on
the original theta complex, retaining the signs on derivative slots.
The resulting calculation also determines precisely what its bare
$L^2(dx)$ completion retains. The Hilbert space and the physical
coordinate remain the original ones throughout.

\begin{theorem}[Reflection on the original complex and every full primary]
On $\B$ define
\[
 (\mathcal IF)(x)=x^{-1}F(1/x),\qquad (\mathcal CF)(x)=\overline{F(x)}.
 \tag{RC1}
\]
Use the source Fourier convention
$\widehat\phi(y)=\int_\R\phi(x)e^{-2\pi ixy}\,dx$.
Then $\widehat V=V$, and
\[
 \mathcal I\Theta\phi=\Theta\widehat\phi,
 \quad\mathcal C\Theta\phi=\Theta\overline\phi,
 \quad D\mathcal I=\mathcal I(1-D).
 \tag{RC2}
\]
The operator $\mathcal I$ is a self-adjoint unitary involution on
$L^2(dx)$, while $\mathcal C$ is an antiunitary involution. They
commute on $\B$, descend to $Q$, and preserve its finite primary
subspace $\sigma_ZE_Z$ whenever $Z$ is closed under reflection and
conjugation, with every original multiplicity retained.

Write a jet as $u_\rho(t)=\sum_{j<m_\rho}u_{\rho,j}t^j$.
The exact induced actions are
\[
 (\mathsf Wu)_{\rho,j}=(-1)^j u_{1-\rho,j},\qquad
 (\mathsf Cu)_{\rho,j}=\overline{u_{\bar\rho,j}}.
 \tag{RC3}
\]
For the original jet map and canonical section,
\[
 J_Z\mathcal I=\mathsf WJ_Z,\quad J_Z\mathcal C=\mathsf CJ_Z,
 \quad \mathcal IR_Z=R_Z\mathsf W,\quad
 \mathcal CR_Z=R_Z\mathsf C.
 \tag{RC4}
\]
In particular, if $f_r=D^rf_0$, then
\[
 \mathcal If_0=f_0,\qquad
 \mathcal If_r=(1-D)^rf_0
 =\sum_{j=0}^r(-1)^j\binom rj f_j.
 \tag{RC5}
\]
\end{theorem}
\begin{proof}
Fourier transformation preserves even Schwartz functions and
interchanges the conditions $\phi(0)=0$ and $\int\phi=0$.
Poisson summation, with these two terms zero, gives (RC2).
The same formula also follows from the fully specified source
construction in \cite[Lemma 4.1]{CCcycles}. The chain rule gives
$D\mathcal I=\mathcal I(1-D)$ directly. Both endpoints defining
$\B$ remain rapidly decreasing under $\mathcal I$ and $\mathcal C$.
The substitution $y=1/x$ proves
$\int|x^{-1}F(1/x)|^2dx=\int|F(y)|^2dy$.
Since $\mathcal I^2=1$, it is self-adjoint; the other Hilbert and
commutation assertions follow directly from the displayed formulas.

In the original Mellin convention,
$\M\mathcal IF(w)=\M F(1-w)$ and
$\M\mathcal CF(w)=\overline{\M F(\bar w)}$.
Taylor expansion at each actual root proves (RC3), including every
factor $(-1)^j$, and the first two identities of (RC4).
For $h=h_Z$ of degree $d$, root stability and multiplicities give
$h(1-w)=(-1)^dh(w)$ and $\overline{h(\bar w)}=h(w)$.
Retain $g=2\xi$, so $g(1-w)=g(w)$ and
$\overline{g(\bar w)}=g(w)$. The original polynomial
$r_Z(u)$ is the unique polynomial of degree less than $d$ representing
$(h/g)u$ modulo $h$. Therefore uniqueness gives
\[
 r_Z(\mathsf Wu)(w)=(-1)^d r_Z(u)(1-w),\qquad
 r_Z(\mathsf Cu)(w)=\overline{r_Z(u)(\bar w)}.
\]
Insert these identities in $\M R_Zu=g r_Z(u)/h$.
Mellin uniqueness proves the last two equalities in (RC4)
as actual function identities, before passage to cohomology.
Finally $\M f_0=g$ and its functional equation prove
$\mathcal If_0=f_0$, after which the operator relation in (RC2)
proves the exact binomial formula (RC5).
\end{proof}

\begin{theorem}[Density of the original theta derivative family]
The original family itself satisfies
\[
 \overline{\operatorname{span}\{f_r:r\ge0\}}^{\,L^2(dx)}
 =L^2(\R_{>0},dx),\qquad
 \overline{\Theta V}^{\,L^2(dx)}=L^2(\R_{>0},dx).
 \tag{RC6}
\]
For every fixed original finite divisor $Z$, the exact attained
quotient metrics
\[
 \begin{split}
 G_M&=(\langle f_i,f_j\rangle)_{0\le i,j\le M},\\
 B_Mu&=(\langle f_i,R_Zu\rangle)_{0\le i\le M},\\
 H_Z^{(M)}&=R_Z^*R_Z-B_M^*G_M^{-1}B_M
 \end{split}
 \tag{RC7}
\]
obey $G_M>0$, $H_Z^{(M)}>0$ at each finite $M$, and
\[
 H_Z^{(M+1)}\preceq H_Z^{(M)},\qquad
 \lim_{M\to\infty}\|H_Z^{(M)}\|=0.
 \tag{RC8}
\]
The inner product is $\langle F,G\rangle=\int\overline F G\,dx$.
Here (RC7) uses precisely the displayed relation subspace
$\operatorname{span}(f_0,\ldots,f_M)$; it does not replace another
programme denominator by that subspace.
\end{theorem}
\begin{proof}
The unitary logarithmic-coordinate map is
$F\mapsto e^{y/2}F(e^y)$ from $L^2(dx)$ to $L^2(dy)$.
Fourier Plancherel therefore identifies the original space with
$L^2(\R,dt/(2\pi))$ by $F\mapsto\M F(1/2+it)$.
This use of Fourier Plancherel is classical; see \cite{SteinFourier}.
Put $g_0(t)=2\xi(1/2+it)$. The images of $f_r$ are
$(1/2+it)^rg_0(t)$, retaining the original generator $D$.

We first prove the exponential integrability actually needed below.
For $s=1/2+it$, the alternating-series continuation obeys
\[
 \eta(s)=s\int_1^\infty A(y)y^{-s-1}\,dy,
 \quad A(y)=\sum_{n\le y}(-1)^{n-1},\quad |A(y)|\le1.
\]
Absolute convergence of this integral gives $|\eta(s)|\le2|s|$.
The identity $\eta(s)=(1-2^{1-s})\zeta(s)$, proved first in
$\Re s>1$ and then continued through $\Re s>0$, gives
$|\zeta(1/2+it)|\le2|1/2+it|/(\sqrt2-1)$.
The vertical-line consequence of Stirling's formula
\cite[5.11.3]{DLMFStirling} gives
\[
 |g_0(t)|\le C(1+|t|)^{11/4}e^{-\pi|t|/4}.
\]
Enlarging $C$ covers the compact interval around zero. Consequently
\[
 \int_\R |g_0(t)|^2e^{2a|t|}\,dt<\infty
 \qquad(0<a<\pi/4).
 \tag{RC9}
\]
The zeros of the nonzero entire function $\xi$ are discrete, so
$g_0$ is nonzero almost everywhere; this uses no zero-location
hypothesis.

Suppose $v\in L^2(dt/(2\pi))$ is orthogonal to all the displayed
polynomial multiples. Then $k(t)=\overline{v(t)}g_0(t)$ belongs
to $L^1(e^{a|t|}dt)$ for each $a<\pi/4$, by Cauchy--Schwarz
and (RC9), and every moment of $k$ is zero. Its bilateral Laplace
transform $K(z)=\int k(t)e^{zt}dt/(2\pi)$ is holomorphic on
$|\Re z|<\pi/4$: on each smaller closed strip its derivatives
are dominated using a slightly larger $a$ in (RC9).
All derivatives at zero vanish, hence the identity theorem gives
$K=0$ on the strip. In particular its Fourier transform vanishes.
The uniqueness theorem for the Fourier transform of an $L^1$
function \cite{SteinFourier} gives $k=0$ almost everywhere.
Since $g_0\ne0$ almost everywhere, $v=0$. This proves density.
Each $f_r=\Theta D^r\phi_*$ belongs to $\Theta V$, proving (RC6).

Linear independence of the $f_r$ follows in the same Mellin
representation from independence of polynomials on the real line,
so $G_M>0$. Orthogonal projection onto their finite span gives
exactly the Schur minimum (RC7). If $H_Z^{(M)}$ vanished on
$u$, then $R_Zu$ would be a linear combination of the $f_r$.
Applying the original full jet map gives $u=0$, since it kills
every $f_r$ and $J_ZR_Z=1$. Thus the finite minimum is positive
definite. Its monotonicity follows from inclusion of the finite
relation spans. Density shows that each diagonal matrix entry
tends to zero; positivity bounds each off-diagonal entry by the
geometric mean of the two diagonal entries. Since the fixed
divisor has finite dimension, this proves operator-norm convergence
in (RC8).
\end{proof}

\begin{proposition}[Exact comparison of the two quotient topologies]
The inclusion of the original $\B$ in $L^2(dx)$ induces the map
\[
 Q=\B/\Theta V\longrightarrow
 L^2(dx)/\overline{\Theta V}^{\,L^2(dx)}=0.
 \tag{RC10}
\]
The infimum of the original $L^2(dx)$ norm in every coset of $Q$
is zero. Nevertheless, for each original finite divisor $Z$,
$\overline J_Z\sigma_Z=1$ on $E_Z$ in the original cohomology.
No nonzero component of $J_Z$ extends to a continuous linear
functional on $L^2(dx)$. The original rapidly decreasing topology
on $\B$ keeps all of these finite jets continuous.
\end{proposition}
\begin{proof}
The target and coset infimum follow from (RC6); the actual
inclusion and quotient maps give (RC10). More explicitly, for a
fixed $u\in E_Z$ take the minimizers
\[
 F_{M,u}=R_Zu-\sum_{r=0}^M (G_M^{-1}B_Mu)_r f_r.
\]
Their original jets equal $u$, their classes in $Q$ all equal
$\sigma_Zu$, and their squared norms equal
$u^*H_Z^{(M)}u\to0$. A nonzero jet functional would therefore
violate continuity if extended to $L^2(dx)$. In the original
$\B$ topology its defining integral has integrand
$F(x)x^{\rho-1}(\log x)^j/j!$; a sufficiently strong rapid-decay
seminorm at each endpoint bounds this integral, proving continuity
there. The finite section identity was proved in (T2).
\end{proof}

This is a calculated receiving map between the algebraic/rapid-decay
cohomology and its bare unweighted Hilbert completion. It does not
close the RH programme. It proves that this particular completion
forgets every zero class, including critical-line classes. The
finite attained metrics (RC7), their original jet maps, and the
original topological quotient specify precisely what remains
available for further arithmetic estimates. No uniform rate in a
growing divisor or in the programme's other quotient flags follows
from the fixed-divisor convergence (RC8).

\section{The full Hurwitz pole train and a bilateral original-jet receiver}
The earlier compact CUE/Hurwitz manuscripts compute their principal
parts at $0$ and $1$ exactly. For a recipient in the programme's
original rapidly decreasing space, all the other gamma poles must
also be retained. The following construction computes them and
then supplies an actual bilateral lattice recipient.

\begin{theorem}[All negative-even residues of the original compact shifts]
Let $\mu$ be the earlier compactly supported probability measure
on $[0,L]$, and retain its actual shift parameter $t\ge0$. Write
\[
 \Lambda_{\mu,t}(w)=\pi^{-w/2}\Gamma(w/2)
             \int\zeta(w,1+ty)\,d\mu(y).
\]
At every integer $k\ge1$ its residue is
\[
 \operatorname{Res}_{w=-2k}\Lambda_{\mu,t}
 =-\frac{2(-1)^k\pi^k}{k!(2k+1)}
            \int B_{2k+1}(1+ty)\,d\mu(y).
 \tag{HU1}
\]
For $t>0$, the residue at $-2$ is strictly positive unless
$\mu=\delta_0$:
\[
 \operatorname{Res}_{w=-2}\Lambda_{\mu,t}
   =\frac{2\pi}{3}\int ty(ty+\tfrac12)(ty+1)\,d\mu(y).
 \tag{HU2}
\]
Consequently removing only the $0,1$ principal parts does not
produce an entire Mellin transform for a nontrivial positive
compact shift. The two-endpoint formula in the earlier paper
continues to hold exactly on its stated two-endpoint scope.
\end{theorem}
\begin{proof}
The classical Hurwitz value is
$\zeta(-n,a)=-B_{n+1}(a)/(n+1)$; see
\cite[Chapter 12]{ApostolHurwitz}, documented in DLMF 25.11.14.
Compactness of $\mu$ permits integration of this identity and of
the holomorphic factors in a neighborhood of $-2k$.
The residue of $\Gamma(w/2)$ there is $2(-1)^k/k!$,
and $\pi^{-w/2}=\pi^k$ at that point. These give (HU1).
The exact polynomial
$B_3(1+a)=a(a+1/2)(a+1)$ is positive for $a>0$ and zero
at $a=0$, proving (HU2) and its equality case.
The $0,1$ principal parts are holomorphic at $-2$, so subtracting
them leaves this nonzero residue unchanged.
\end{proof}

\begin{theorem}[Bilateral completion in the original physical coordinate]
For $0<a<1$ define
\[
 \begin{split}
 K_a(x)&=2\sum_{n\in\mathbb Z}e^{-\pi(n+a)^2x^2},\\
 F_a(x)&=D(D-1)K_a(x),\\
 A(w)&=w(w-1)\pi^{-w/2}\Gamma(w/2).
 \end{split}
 \tag{HU3}
\]
For each closed interval $[\epsilon,1-\epsilon]$ with
$0<\epsilon<1/2$, $F_a$ belongs to the original $\B$ uniformly
in every defining endpoint seminorm. Its entire Mellin transform is
\[
 \M F_a(w)=A(w)\bigl(\zeta(w,a)+\zeta(w,1-a)\bigr)
 =A(w)\bigl(a^{-w}+\zeta(w,1+a)+\zeta(w,1-a)\bigr).
 \tag{HU4}
\]
Thus the two earlier shifts together with the actual missing
central summand cancel every negative-even gamma pole. The exact
physical inversion of this family is
\[
 \mathcal IF_a(x)=4\sum_{n\ge1}
                \cos(2\pi na)\,\phi_*(nx).
 \tag{HU5}
\]
This retains the full original gaussian source $\phi_*$, the
original coordinate $x$, and every integer lattice coefficient.
\end{theorem}
\begin{proof}
For $a$ in the stated interval, every $|n+a|$ is bounded away
from zero, and the gaussian sum and all its $D$ derivatives
decrease exponentially as $x\to\infty$, uniformly in $a$.
Poisson summation with the Fourier convention in (RC2) gives
\[
 K_a(x)=\frac2x\sum_{m\in\mathbb Z}
             e^{-\pi m^2/x^2}e^{2\pi ima}.
 \tag{HU6}
\]
The term $m=0$ is exactly $2/x$ and is killed by $D(D-1)$.
All other terms, with every $D$ derivative, decrease faster than
every power at zero uniformly in $a$. This proves membership in
the original $\B$ at both ends, not only a formal pole cancellation.

In $\Re w>1$, absolute integration of the gaussian series gives
\[
 \M K_a(w)=\pi^{-w/2}\Gamma(w/2)
                  \bigl(\zeta(w,a)+\zeta(w,1-a)\bigr).
\]
Indeed positive terms $n+a$ for $n\ge0$ give the first Hurwitz
series and absolute values of terms with $n\le-1$ give the second.
The factor $2$ in $K_a$ cancels the $1/2$ in each gaussian
Mellin integral. Two integrations by parts are justified in this
half-plane at both endpoints, yielding the first equality of
(HU4). Its left-hand side is entire by the just-proved endpoint
estimates. The classical recurrence
$\zeta(w,a)=a^{-w}+\zeta(w,1+a)$ proves the second equality.
At negative even integers, the Bernoulli reflection identity
$B_{2k+1}(1-a)=-B_{2k+1}(a)$ gives the same cancellation
algebraically; it follows also directly from the gaussian proof.

Apply $\mathcal I$ to (HU6). It gives
$\mathcal IK_a=2+4\sum_{n\ge1}\cos(2\pi na)e^{-\pi n^2x^2}$.
The identity $D(D-1)\mathcal I=\mathcal ID(D-1)$ follows
from (RC2). This polynomial kills the constant and maps each
gaussian to the actual $\phi_*(nx)$, proving (HU5).
All termwise differentiations are absolutely locally convergent.
Poisson summation and the gaussian transform used here are
classical, with their Fourier convention fixed in
\cite{SteinFourier}; no new convention for $x$ or for the Mellin
coordinate has been introduced.
\end{proof}

\begin{theorem}[Every original finite primary is attained by this family]
Let $Z$ be any nonempty finite set of actual nontrivial zeros,
with their full multiplicities, and let $J_Z:\B\to E_Z$ be the
original jet map in (T1). For every $0<\epsilon<1/2$ the
vectors
\[
 j(a):=J_ZF_a\qquad(\epsilon<a<1-\epsilon)
\]
span all of $E_Z$. In the original ordered jet coordinates set
\[
 \mathsf P_{Z,\epsilon}
    =\int_\epsilon^{1-\epsilon} j(a)j(a)^*\,da.
 \tag{HU7}
\]
Then $\mathsf P_{Z,\epsilon}>0$, and the explicit map
\[
 \mathsf S_{Z,\epsilon}u
   =\int_\epsilon^{1-\epsilon}
      F_a\,j(a)^*\mathsf P_{Z,\epsilon}^{-1}u\,da
       \quad:E_Z\longrightarrow\B
 \tag{HU8}
\]
satisfies
\[
 J_Z\mathsf S_{Z,\epsilon}=1,\qquad
 J_Z(\mathsf S_{Z,\epsilon}-R_Z)=0.
 \tag{HU9}
\]
The actual original finite-sector projection on its cohomology
classes is
$\sigma_Z\overline J_Z[q\mathsf S_{Z,\epsilon}u]=\sigma_Zu$.
The residual difference remains in the explicitly stated
$\ker J_Z$; it is not identified with $\Theta V$.
\end{theorem}
\begin{proof}
The components of $j(a)$ are real-analytic functions on $0<a<1$,
by (HU4) and the holomorphy of the Hurwitz parameter there.
Suppose a complex linear functional $\lambda$ on $E_Z$ kills
them on the stated open interval. Real-analytic uniqueness extends
this identity to all $0<a<1$.

Near $a=0$, the sum
$\zeta(w,1+a)+\zeta(w,1-a)$ and all its finitely many required
$w$ derivatives are holomorphic and bounded in $a$.
The singular part of $\lambda j(a)$ therefore has the exact form
\[
 \sum_{\rho\in Z}a^{-\rho}P_\rho(-\log a),
       \qquad \deg P_\rho<m_\rho.
 \tag{HU10}
\]
The original jet convention is the coefficient of $(w-\rho)^j$,
so these polynomials include their original $1/j!$ factors.
For each root the map from $\lambda$'s primary coefficients to
the coefficients of $P_\rho$ is triangular, with nonzero diagonal
multipliers $A(\rho)/j!$. They are nonzero since
$0<\Re\rho<1$, $\rho\ne0,1$, and gamma has no zeros.
Thus a nonzero $\lambda$ gives at least one nonzero polynomial.

Such an exponential polynomial cannot stay bounded as $a\to0$.
For completeness put $y=-\log a$, choose the largest real part
$\beta>0$ among its nonzero terms, and then the largest degree $d$
at that real part. After division by $e^{\beta y}y^d$, it equals
$T(y)+o(1)$, uniformly as $y\to\infty$, where
$T(y)=\sum c_\rho e^{i\Im\rho y}$ is a nonzero finite sum with
distinct frequencies. If the original expression were bounded,
then $T(y)\to0$. On the other hand direct integration gives
\[
 \lim_{R\to\infty}\frac1R\int_R^{2R}|T(y)|^2dy
       =\sum |c_\rho|^2>0:
\]
each cross-frequency integral divided by $R$ tends to zero,
while each diagonal integral is $|c_\rho|^2$. This contradicts
$T(y)\to0$. Since the remainder of $\lambda j(a)$ is bounded,
the identity $\lambda j=0$ consequently forces $\lambda=0$.
This proves the spanning assertion without discarding any primary
or derivative slot.

For $u\ne0$, the continuous nonnegative integrand
$|j(a)^*u|^2$ cannot vanish on the whole interval by the spanning
result. Its integral is therefore positive, proving (HU7).
The uniform endpoint seminorm bounds in (HU3) show that (HU8)
is a genuine integral in $\B$. Its continuous jet map may be
passed through the integral, giving
$J_Z\mathsf S_{Z,\epsilon}u
=\mathsf P_{Z,\epsilon}\mathsf P_{Z,\epsilon}^{-1}u=u$.
Subtracting the already proved $J_ZR_Z=1$ proves (HU9).
Passage through the original quotient $q$ proves the stated
finite-sector cohomology formula. This calculation supplies the
actual map and its exact residual space, without assigning a
false equality to the two sections in the entire quotient.
\end{proof}

The finite guard $\epsilon$ is retained: no uniform statement at
$a=0$ is asserted. The interpolation matrix (HU7) is explicitly
evaluated as an integral of full Hurwitz jets, not as a rank-only
allocation or an arithmetic sign estimate. It gives an additional
concrete test-family receiver from the older CUE/Hurwitz construction
to every original finite zero sector. Its residual cohomology class
and its original $dx$ Gram matrix remain specified objects for
subsequent arithmetic work.

\begin{thebibliography}{99}
\bibitem{ApostolHurwitz} T. M. Apostol, \emph{Introduction to Analytic
Number Theory}, Springer, 1976, Chapter 12; the Hurwitz recurrence
and Bernoulli special values are documented with primary source
locators in
\href{https://dlmf.nist.gov/25.11.E3}{NIST DLMF 25.11.3} and
\href{https://dlmf.nist.gov/25.11.E14}{25.11.14}.
\bibitem{DLMFStirling} Richard A.~Askey and Ranjan Roy,
NIST Digital Library of Mathematical Functions,
Section 5.11, Asymptotic expansions of the gamma function,
\href{https://dlmf.nist.gov/5.11.E3}{formula 5.11.3}, including the
classical Stirling expansion. Its vertical-line modulus estimate
is used in (RC9).
\bibitem{SteinFourier} E. M. Stein and R. Shakarchi,
\emph{Fourier Analysis: An Introduction}, Princeton University
Press, 2003, Chapter 5 (Fourier transformation on the real line,
inversion and Plancherel).
\bibitem{ApostolMobius} T. M. Apostol, \emph{Introduction to Analytic
Number Theory}, Springer, 1976, Chapter 2; the divisor identity and
dilated-sum inversion are documented with this source in
\href{https://dlmf.nist.gov/27.5.E2}{NIST DLMF 27.5.2} and
\href{https://dlmf.nist.gov/27.5.E7}{27.5.7}.
\bibitem{StacksSupport} The Stacks Project Authors, \emph{The Stacks Project},
Section 20.21, Cohomology with support in a closed subset,
\href{https://stacks.math.columbia.edu/tag/0A39}{Tag 0A39}.
The defining derived support functor used in (S2) is standard;
(S4) supplies the full injective-resolution calculation for the exact
two-point space and original theta arrow used here.
\bibitem{Beurling} A. Beurling, A closure problem related to the Riemann
zeta-function, \emph{Proc. Natl. Acad. Sci. USA} 41 (1955), 312--314,
\href{https://doi.org/10.1073/pnas.41.5.312}{doi:10.1073/pnas.41.5.312}.
\bibitem{BD} L. B\'aez-Duarte, A strengthening of the Nyman--Beurling
criterion for the Riemann hypothesis, \emph{Rend. Lincei Mat. Appl.}
14 (2003), 5--11; \href{https://arxiv.org/abs/math/0202141}{arXiv:math/0202141}.
\bibitem{CC} A. Connes and C. Consani, Geometry of the arithmetic site,
\emph{Adv. Math.} 291 (2016), 274--329,
\href{https://doi.org/10.1016/j.aim.2015.11.045}{doi:10.1016/j.aim.2015.11.045}.
The framed denominator construction cited above is in the earlier local
programme manuscript; it must not be attributed to this human paper
without a matching source statement.
\bibitem{DLMF25} T. M. Apostol, \emph{Introduction to Analytic Number Theory},
Springer, 1976, Section 12.8; formulas also documented in
\href{https://dlmf.nist.gov/25.4.E4}{NIST DLMF 25.4.4}.
\bibitem{CCcycles} A. Connes and C. Consani, Hochschild homology,
trace map and $\zeta$-cycles, \emph{Proc. Sympos. Pure Math.} 105
(2023), 83--101, especially Lemma 4.1 and Proposition 4.2;
\href{https://arxiv.org/abs/2207.10419}{arXiv:2207.10419}.
\bibitem{Titchmarsh} E. C. Titchmarsh, \emph{The Theory of the Riemann
Zeta-Function}, second edition, revised by D. R. Heath-Brown,
Oxford University Press, 1986 (the classical zero-counting theorem used
in the nonzero-tail assertion).
\bibitem{DeligneII} P. Deligne, La conjecture de Weil: II,
\emph{Publications Math\'ematiques de l'IH\'ES} 52 (1980), 137--252,
especially (1.6.1), (1.6.7)--(1.6.10);
\href{https://www.numdam.org/item/PMIHES_1980__52__137_0/}{primary article},
doi:10.1007/BF02684780. The finite primary formulas and every metric
receiver used here are proved explicitly above.
\bibitem{Green} P. Green, The local structure of twisted covariance
algebras, \emph{Acta Mathematica} 140 (1978), 191--250,
doi:10.1007/BF02392308;
\href{https://archive.ymsc.tsinghua.edu.cn/pacm_download/117/6237-11511_2006_Article_BF02392308.pdf}{primary article}.
\end{thebibliography}
\end{document}
